% \begin{savequote}
% \includegraphics[width=5cm]{./images/teaser/follow-me-kontur.jpg}
% \end{savequote}
\def\ChapterCode{PAM}
	\chapter
[\CO{[PAM]} Patientenmanagement und standardisierte Patientenversorgung]
{Patientenmanagement und standardisierte Patientenversorgung}
\label{chp:PAM}

\ChapterIntro
{%
Unter dem Begriff Patientenmanagement fasst man alle
Maßnahmen und Entscheidungen zusammen, die der Untersuchung,
Betreuung und Behandlung von Patienten und Abhandlung von
Einsätzen und anderen patientenbezogenen Aufträgen dienen,
sowie die richtige und angemessene zeitliche Abfolge dieser
Maßnahmen.

Dieses Kapitel behandelt die richtige Kombination der in
anderen Kapiteln beschriebenen Maßnahmen in einen
funktionierenden und vorteilhaften Ablauf der
Patientenversorgung. Grundbausteine sind dabei die Kenntnis
und das Beherrschen von medizinischem \EE{Wissen} und
sanitätshilflichen \EE{Fertigkeiten}. Das Erarbeiten von
(Verdachts-) Diagnosen mittels \EE{Patientengespräch} und
\EE{Untersuchungen} sowie das Beherrschen der
\EE{theoretisch-medizinischen Grundlagen} sind weitere
Voraussetzungen für effektives und professionelles Arbeiten.
}
{%
\begin{description}
  \item[Maintainer:] Sebastian Gabriel
  \item[Autoren:] Diverse
  \item[Reviewer:] Standard-Reviewprozess
  \item[Version:] {\VarDocVersionTypeName} ({\VarDocVersionTypeDescription})
  \item[SHA1:] {\DocGitCommit}
\end{description}

{\TxTeilkapitelAASS}}

% \includegraphics[width=\linewidth]{./images/emhofer-josef-aass/Anamnese_Symbolbild_75pc.jpg}


% \the\textwidth
% 
% \the\marginparsep
% 
% \the\marginparwidth

\clearpage % TEMP temp clearpage
% \TextSlide{\TxQuerverweise}
% {~
% 
% \noindent\includegraphics[angle=270,width=\linewidth]{./images/leitsystem/Leitschema-PAM_gesamt.pdf}
% }{\begin{itemize}
%   \item \EE{Medizinprodukte allgemein}: \CO{[MP1]} (\ref{chp:MP1})
%   \item \EE{Untersuchungen}:                 \CO{[BAU]} (\ref{chp:BAU})
%   \item \EE{Sanitätshilfemaßnahmen}:         \CO{[SAN]} (\ref{chp:SAN})
%   \item \EE{Anamnese und Patientengespräch}: \CO{[AUP]} (\ref{chp:AUP})
% \end{itemize}}{}{}{}{}

% \section{Grundsätzlicher Ablauf einer präklinischen Hilfeleistung}



% \begin{WidePar}
\pagestyle{empty}\loadgeometry{WholePage}
\vfill

\begin{minipage}{\linewidth}
\noindent{\sffamily\color{Base}%
{\bfseries\Huge\color{Base}
Standardisierte Patientenversorgung}\bigskip 

\color{BaseMiddle}Arbeits- und Ausbildungsstandards für den \color{DefaultRed}Sanitätsdienst} | {\AASS} \opt{VAR-VAR2}{| {\color{AsboeRed}ASBÖ}}
\hfill\textsf{\VarDocVersion}

\bigskip

\Layout{57}{}{}{}
{%
\begin{center}
  \includegraphics[width=280pt]{./images/emhofer-josef-aass/PAMMaerz2014Vers2}
\end{center}
}
{Graphische Übersicht: \EEE{Standardisierte Patientenversorgung}}
{}



% \vspace{2\baselineskip}
\vfill

\Layout{57}{}{}{}{%
{\sffamily\begin{tabularx}{\linewidth}{XXX}
\addlinespace\toprule\addlinespace
\noindent\TabHeading{Massive Störung}
&
\noindent\TabHeading{Alarmsymptome}
&
\noindent\TabHeading{Alarmdiagnosen}
\\\addlinespace\midrule\addlinespace
Bewusstsein

Atemweg / Atmung 

Kreislauf 

(\ABCDE{3}\ABCDE{A}\ABCDE{B}\ABCDE{C})
&
Atemnot, die sich nicht bessert 

Brodelndes Atemgeräusch 

Thoraxschmerz  

Schocksymptome 

entgleiste Vitalwerte

schwere Verletzungen

Hirndruckzeichen

{\dots} 
&
Herzinfarkt

Kardiales Lungenödem

Status Epilepticus

Beckenfraktur

Rückenmarksverletzung

\ldots
\\\addlinespace\bottomrule
\end{tabularx}}
}{{\color{DefaultRed}\bfseries{\VarFlag}} \Ca{Alarmzeichen:}
Wann ist eine vitale Bedrohung
wahrscheinlich? Eine Übersicht.\label{tab:alarmzeichen}}{}

\vfill
\noindent{\sffamily\small{\textcopyright} \today\ \ARGEAASS \hfill \url{http://www.aass.at}}
\end{minipage}
\clearpage
% 
% \vspace{-\baselineskip}~
~
\vfill



% \MassnahmenWide{Standardmaßnahmen}{Einschätzungsblock}



\begin{tcolorbox}[
title={Standardmaßnahme YY12100B: Einschätzungsblock
% \captionof{massnahmen}[#1: \E{#2}]{\color{Base}\E{#1}: #2 #3\hrulefill{\fontfamily{lmr}\selectfont\textrecipe}}
},
fonttitle=\sffamily\bfseries\small,fontupper=\sffamily\small,
left=2pt,right=2pt,
colframe=ProceduresBoxBorder,
colback=MarginBackground,
% skin=beamer,
]
\renewcommand{\labelitemi}{\textcolor{LabelItemBoxProcedures}{\SymbolLabelItemI}}%
\renewcommand{\labelitemii}{\textcolor{LabelItemBoxProcedures}{\small\SymbolLabelItemII}}%
\flushleft{}\CorrectionListsBox
{\label{m:ersteinschaetzung}\label{m:einschaetzungsblock}}

\CorrectionItemizeBox\renewcommand{\itemhook}{%
%   \setlength{\topsep}{0pt}%
  \setlength{\itemsep}{0.004\baselineskip}
%   \setlength{\parskip}{\baselineskip}
  \setlength{\leftmargin}{1.4em}
  \setlength{\labelwidth}{0.9em}
}
\noindent\begin{minipage}[t]{\linewidth -
{2\fboxrule} - {2\fboxsep}}
\noindent\begin{minipage}[t]{\linewidth}
{%\vspace{-2em}%
\footnotesize
% \begin{tabulary}{\linewidth}[H]{cp{8em}p{8.5cm}}
\begin{tabularx}{\linewidth}[H]{p{2cm}XX}
% \toprule\addlinespace
&
\TabHeading{Beurteilung / Untersuchung}
&
\TabHeading{Typische Sofortmaßnahmen}
\\\addlinespace\toprule\addlinespace
\noindent\TabHeading{\ABCDEPlain{1}}\vspace{\baselineskip}

\TabHeading{Szeneüberblick m. Schutz}
&
\begin{minipage}[t]{\linewidth}
\begin{itemize}
  \item Gefahrenzonen, Sicherheit, Selbstschutz;
  \item Umstände, Umgebung, Ort,
  Zeit, Patientenanzahl;
  \item \EE{Trauma?} \EE{Unfallmechanismus?}
  \item Weitere Ressourcen erforderlich?,
  Großschaden?
  \item Lagemeldung erforderlich
\end{itemize}
\end{minipage}
&
\noindent\color{DefaultA}\begin{minipage}[t]{\linewidth}
\begin{itemize}
  \item GAS-Maßnahmen (\prettyref{topic:gefahrengutunfall})
  \item Lagemeldung
  \item Schutzausrüstung
  \item Nachforderung weiterer Kräfte
\end{itemize}
\end{minipage}
\\\addlinespace\toprule\addlinespace
\noindent\TabHeading{\ABCDEPlain{2}}\vspace{\baselineskip}

\TabHeading{Eindruck}
&
\begin{minipage}[t]{\linewidth}
\begin{itemize}
  \item Alter, Geschlecht, \EE{Hautfarbe}, Gesichtsausdruck,
  \EE{Haltung}, spontane Bewegungen,  Sprache
  \item Offensichtliche Verletzungen, Blutungen
  \item Sonstige \EE{Auffälligkeiten}
\end{itemize}
\end{minipage}
&
\noindent\color{DefaultA}\begin{minipage}[t]{\linewidth}
\begin{itemize}
  \item Manuelle HWS-Fixierung
  \item Stillung von starken Blutungen
  \item Bewegungsverbot
\end{itemize}
\end{minipage}
\\\addlinespace\midrule\addlinespace
\noindent\TabHeading{\ABCDEPlain{3}}

\TabHeading{Bewusstsein}
&
\begin{minipage}[t]{\linewidth}
\begin{itemize}
  \item \EE{Bewusstseinsgrad} (WASB)
\end{itemize}
\end{minipage}
&
\noindent\color{DefaultA}\begin{minipage}[t]{\linewidth}
\begin{itemize}
  \item NA-Nachforderung
\end{itemize}
\end{minipage}
\\\addlinespace\midrule\addlinespace
\noindent\TabHeading{\ABCDEPlain{4}}

\TabHeading{Hauptbeschwerde}
&
\begin{minipage}[t]{\linewidth}
\begin{itemize}
  \item \EE{Berufungsgrund} und
  \EE{Leitsymptome}
\end{itemize}
\end{minipage}
&
\noindent\color{DefaultA}\begin{minipage}[t]{\linewidth}
\end{minipage}
\\\addlinespace\toprule\addlinespace
\noindent\TabHeading{\ABCDEPlain{A}}\vspace{\baselineskip}

\TabHeading{Atemweg \newline Airway}
&
\begin{minipage}[t]{\linewidth}
\begin{itemize}
  \item *Inspektion der Atemwege (\EE{Mund}, \EE{Nase})
  \item \EE{Atemgeräusche}
\end{itemize}
\end{minipage}
&
\noindent\color{DefaultA}\begin{minipage}[t]{\linewidth}
\begin{itemize}
  \item Absaugung, Fremdkörper entfernen
  \item Kopf überstrecken, Esmarch-Handgriff
%   \item HWS-Immobilisation manuell oder mit Schiene
  \item \Z{Erweiterte Maßnahmen}
\end{itemize}
\end{minipage}
\\\addlinespace\midrule\addlinespace
\noindent\TabHeading{\ABCDEPlain{B}}\vspace{\baselineskip}

\TabHeading{Atmung\newline (Breathing)}
&
\begin{minipage}[t]{\linewidth}
\begin{itemize}
  \item Schätzen: \EE{Atemfrequenz}, \EE{-tiefe}
  \item *\EE{Sp\ce{O2}}
  \item Atemgeräusche
  \item Inspektion der \EE{Thorax-Atembewegungen}
  \item \E{Zeichen der
  Atemnot}: Hautfarbe, Bauch-/Thorax\-bewegungen,
  Anstrengung beim Atmen, Atemhilfsmuskulatur, \ldots
  \item \Z{(Auskultation: Seitenvergleich, Lungenbasen)}
\end{itemize}
\end{minipage}
&
\noindent\color{DefaultA}\begin{minipage}[t]{\linewidth}
\begin{itemize}
  \item Reanimation
  \item Assistierte Beatmung (bei AF \textless\,8 oder
  \textgreater\,30\PerMin*)
  \item \ce{O2}-Gabe
  \item Situationsgerechte Lagerung
  \item \Z{Erweiterte Maßnahmen\\(Entlastung
  Spannungspneumothorax, \ldots )}
\end{itemize}
\end{minipage}
\\\addlinespace\midrule\addlinespace
\noindent\TabHeading{\ABCDEPlain{C}}\vspace{\baselineskip}

\TabHeading{Kreislauf \newline (Circulation)}
&
\begin{minipage}[t]{\linewidth}
\begin{itemize}
  \item \EE{Hautfarbe}, \EE{-temperatur}, Schweiß
  \item \EE{Rekap}-Zeit
  \item \EE{Radialispuls}: Stärke, Rhythmus, Abschätzen der
  Frequenz
  \item \EE{Herzfrequenz}
  \item \EE{Blutdruck}
  
\end{itemize}
\hrule
\begin{itemize}
  \item \TabHeading{*Schnelle Trauma-Untersuchung}: Inspektion
  und Abtasten von
  \begin{compactenum}
    \item \EE{Kopf} inkl. Ohren
    \item \EE{Hals}
    \item \EE{Thorax} (Atemprobe)
    \item \EE{Bauch}
    \item \EE{Becken}
    \item \EE{Oberschenkel}
    \item \EE{Rücken} (evtl. erst beim Umlagern)
  \end{compactenum}
\end{itemize}
\end{minipage}
&
\noindent\color{DefaultA}\itshape\begin{minipage}[t]{\linewidth}
\begin{itemize}
  \item Blutstillung
  \item Situationsgerechte Lagerung
\end{itemize}
\end{minipage}
\\\addlinespace\midrule\addlinespace
\noindent\TabHeading{\ABCDEPlain{D}}\vspace{\baselineskip}

\TabHeading{Neurologisch\newline (Disability)}
&
\begin{minipage}[t]{\linewidth}
\begin{itemize}
    \item \EE{Orientierung} oder *\EE{GCS}
    \item \EE{Pupillen}
    \item \EE{Kraftprobe OE}, *UE
    \item *Kann der
    Patient \EE{Hände} und \EE{Füße} spüren und aktiv
    bewegen?
  \item *\EE{Blutzuckermessung}
\end{itemize}
\end{minipage}
&
\noindent\color{DefaultA}\begin{minipage}[t]{\linewidth}
\begin{itemize}
  \item Strategie- und Zielentscheidung (eilig, Rendez-vous, \ldots)
\end{itemize}
\end{minipage}
\\\addlinespace\midrule\addlinespace
\noindent\TabHeading{\ABCDEPlain{E}}

\TabHeading{Erweiterte\newline Untersuchung}

{\tiny (Wenn Einschätzung der vitalen
Bedrohung oder Hauptbeschwerde unklar ist.)}
&
\begin{minipage}[t]{\linewidth}
\begin{itemize}
  \item (Fremd-) \EE{Anamnese} (SAMPLER, OPQRST)
  \item Einschätzen der \EE{Umgebung}
  \item Weitere relevante \EE{Untersuchungen}
  \item Arbeitsdiagnose erstellen und Differentialdiagnosen
  ausschließen
\end{itemize}
\end{minipage}
&
\noindent\color{DefaultA}\begin{minipage}[t]{\linewidth}
\begin{itemize}
  \item Strategie-, Transport- und Zielentscheidung 
  \item Spezielle Maßnahmen gemäß Verdachtsdiagnose
\end{itemize}
\end{minipage}
% \\\addlinespace\bottomrule\addlinespace
% \TabHeading{Beurteilung}
% &
% \begin{minipage}[t]{\linewidth}
% \begin{itemize}
%   \item \EE{Einschätzen der vitalen Bedrohung}
%   ständig während des gesamten Blocks
%   \item Vorläufige Verdachtsdiagnose formulieren
% \end{itemize}
% \end{minipage}
% &
% \noindent\color{DefaultA}\begin{minipage}[t]{\linewidth}
% \begin{itemize}
%   \item {\TxMassVit}
%   \item Reihenfolge der weiteren Untersuchungen und
%   Maßnahmen festlegen
% \end{itemize}
% \end{minipage}
\\\addlinespace\bottomrule
\end{tabularx}

\begin{center}
  \EE{Der Einschätzungsblock muss in regelmäßigen Abständen
in angemessenem Umfang wiederholt
werden (Verlaufskontrolle)!} \tiny

{ Mit einem * versehene Punkte werden nur durchgeführt,
wenn sie \E{situationsentsprechend} sind. \enquote{Typische
Sofortmaßnahmen} sind als \E{Beispiele} zu verstehen. 

Bei \Ca{absolut zeitkritischen Patienten}, bei denen bereits die 
ABC-Einschätzung ergibt dass ein Transport nicht aufschiebbar ist, 
kann u.\,U. schon \Ca{nach dem Punkt C eine vorgezogene Transportentscheidung} 
getroffen werden. 
Die Maßnahmen von D und E sollen dann, sofern möglich, während des 
Transportes erfolgen. }
\end{center}
}
\end{minipage}
\end{minipage}%\vspace{-\baselineskip}
\end{tcolorbox}






\vfill
\restoregeometry\pagestyle{fancy}

% \pagestyle{scrheadings}

% \clearpage




    	







%%%%%%%%%%%%%%%%%%%%%%%%%%%%%%%%%%%%%%%%%%%%%%%%%%%%%%%%%%%

\Layout{57}{}
{}
{}
{%
\begin{minipage}{\linewidth}
\FontTableBase
{\sffamily\FontTableBase\CorrectionItemizeBox\renewcommand{\itemhook}{%
%   \setlength{\topsep}{0pt}%
  \setlength{\itemsep}{0.005\baselineskip}
%   \setlength{\parskip}{\baselineskip}
  \setlength{\leftmargin}{1.4em}
  \setlength{\labelwidth}{0.9em}
}
\begin{tabulary}{\linewidth}[H]{lCL}
\toprule\addlinespace
{\multirow{6}{2cm}{\TabHeading{Anamnese}}}
 &
{\SAMPLER{S}}
 &
Symptome {\&} Schmerzen: OPQRST

{\tiny(Lokalisation, Stärke,  Qualität, Zeit/Verlauf, 
 Ausstrahlung, Linderung/Verstärkung)}
\\\addlinespace
 &
{\SAMPLER{A}}
 &
Allergien/Allergische Reaktion (auch
Allergien auf Medikamente)
\\\addlinespace
 &
 {\SAMPLER{M}}
 &
Medikamente {\tiny(Nimmt Patient Medikamente?
Wogegen sind diese?)}
\\\addlinespace
 &
{\SAMPLER{P}}
 &
Patientengeschichte: Krankheiten {\tiny(chronische, frühere
K.)}, Operationen, \ldots 
\\\addlinespace
 &
 {\SAMPLER{L}}
 &
Letzte \ldots (z.\,B. letzte Mahlzeit, letzte
Regelblutung, letzter Spitalsaufenthalt)
\\\addlinespace
 &
 {\SAMPLER{E}}
 &
Ereignisse vor Notfalleintritt {\tiny(z.\,B.
Anstrengung vor Brustschmerz, Sturz, \ldots)}
\\\addlinespace
 &
 {\SAMPLER{R}}
 &
Risikofaktoren {\tiny(z.\,B.
Herz-Kreislauferkraknungen, Rauchen, Übergewicht, \ldots)}
\\\addlinespace\midrule\addlinespace

{\multirow{8}{2cm}{\TabHeading{Diagnostik}}}
 &
{\sffamily AF}
 &
Atemfrequenz auszählen
\\\addlinespace
 &
({\sffamily Sp\ce{O2}})
 &
Sauerstoffsättigung im Blut, Monitoring
\\\addlinespace
 &
{\sffamily  BZ}
 &
Blutzucker {\tiny(Wann war die letzte
Mahlzeit/Insulinspritze? Diabetes bekannt?)}
\\\addlinespace
 &
{\sffamily Temp.}
 &
Körpertemperatur/Fieber
\\\addlinespace
 &
 {\sffamily\small Neuro\-status}
 &
(Bewusstseinsgrad), Orientierung,
Pupillen, Halbseitenzeichen, Herdblick, Meningismus,
retrograde Amnesie
\\\addlinespace
 &
{\sffamily\small Trauma\-check}
 &
Abnorme Beweglichkeit, DMS an allen
Extremitäten (Nagelbettprobe), Wunden, paradoxe Atmung,
Abwehrspannung des Bauches, ...
\\\addlinespace
 &
{\sffamily\small Körp. Unt.}
&
Sonstige körperliche Untersuchungen: Abdomen abtasten,
Körperstellen inspizieren, \ldots
\\\addlinespace
 &
{\sffamily\small \ldots}
&
EKG, Kapnometrie, \ldots
\\\addlinespace\midrule\addlinespace
\TabHeading{Maßnahmen}
&
&
\begin{minipage}{\linewidth}
\begin{itemize}
  \item \EE{Immer durchzuführende Standardmaßnahmen}
  (\prettyref{m:standardmassnahmen-immer})
  \item Regelmäßige \EE{Wiederholung des
  Einschätzungsblockes} {\tiny(in angemessenem Umfang)}
  \item (Lagerung)
  \item Verband, Blutstillung
  \item Schienung
  \item psychischer Beistand
  \item \EE{Wärmeerhaltung / Kühlung}
  \item Entscheidung Notarzt-Nachforderung wegen
  Schmerztherapie oder Aufklärung/Belassung
  \item Sauerstoffgabe (\FullPrettyRef{m:sauerstoffberieselung})
  \item \ldots
\end{itemize}
\end{minipage}
\\\addlinespace\midrule\addlinespace
\TabHeading{Assistenz}
&
&
\begin{minipage}{\linewidth}
\begin{itemize}
  \item EKG anlegen
  \item Infusion/periphervenösen Zugang vorbereiten
  \item Intubation vorbereiten
  \item Medikamente vorbereiten
  \item \ldots
\end{itemize}
\end{minipage}
\\\addlinespace\midrule\addlinespace
\TabHeading{Transport}
&
&
\begin{minipage}{\linewidth}
\begin{itemize}
  \item Welches Spital/Bett buche ich ab? {\tiny(Arbeitsunfall?
  Welche Abteilung fahre ich an? Ist der Patient bereits in
  einem Spital in Behandlung? Spitalsbehandlung vor kurzer
  Zeit (\enquote{Reklamation})?)}
  
  Brauche ich eine Spezialabteilung? {\tiny(Schockraum, Stroke
  Unit, PTCA, Verbrennungsstation)}
  \item Übergabe {\tiny(Anamnese, Befunde, Maßnahmen, Hinweise)}
\end{itemize}
\end{minipage}

\\\addlinespace\bottomrule
\end{tabulary}}

{\scriptsize\EE{Fett} gedruckte Punkte haben eine besonders hohe Priorität, (eingeklammerte) Punkte sollten schon als Teil des Einschätzungsblockes abgehandelt worden sein.}
\end{minipage}
}
{Kurzübersicht: \EEE{Weiterführende Untersuchungen und Anamnese}.
Die Reihenfolge ist je nach Patient
unterschiedlich!\label{tab:checkliste-patientenmanagement}\label{tab:pam-weiterfuehrende}.}
{}

\clearpage

\section{Allgemeines}
\subsection{Einführung} 
Gerade zu Beginn der Tätigkeit als medizinische
Fachkraft fühlt man sich bei Konfrontation mit Patienten durch die
große Menge an Sinneseindrücken und Informationen überfordert. 
Das Erkennen, was in einer Situation wichtig und was weniger
wichtig ist, ist selbst für erfahrenes Personal oft nicht einfach.
Um diesem Umstand zu begegnen, haben sich strukturierte
Vorgehensweisen bewährt, welche als grundsätzliche Leitlinie
dienen, und sowohl dem Personal, als auch dem Patienten ein
hohes Maß an Sicherheit geben.
Patientenmanagement besteht aus vielen einzelnen
Komponenten, welche wie eine Maschine zusammenarbeiten:
\begin{itemize}
  \item Einschätzungsblock (ABCDE) zum
  \begin{itemize}
    \item Einschätzen und
    Erkennen einer vitalen Bedrohung und
    \item Ergreifen von Sofortmaßnahmen
    \item Erkennen der Hauptbeschwerde und von Leitsymptomen
  \end{itemize}
  \item Weiterführende Untersuchungen und Maßnahmen
  \begin{itemize}
    \item Erhebung der Anamnese (SAMPLER)
    \item Weitere Untersuchungen
    \item Spezielle Maßnahmen je nach Krankheitsbild
  \end{itemize}
  \item Immer durchzuführende Standardmaßnahmen
  (\prettyref{m:standardmassnahmen-immer})
  \item Ggfs. Auswahl einer geeigneten Einrichtung zur
  Weiterbehandlung
  \item Ggfs. Transport
  \item Ggfs. Übergabe an das weiterbehandelnde Personal
\end{itemize}


\Dictum
[Theaterdirektor. \\
In: Johann Wolfgang von Goethe: Faust. Der Tragödie erster Teil]
{Der Worte sind genug gewechselt,\\
Laßt mich auch endlich Taten sehn!}






\clearpage
\section{Einschätzungsblock:
1\,--\,2\,--\,3\,--\,4\,--\,A\,--\,B\,--\,C\,--\,D\,--\,E}
\label{topic:ersteinschaetzung}
\label{topic:einschaetzungsblock}

\begin{quote}
  \Speech{\enquote{Von 1 bis E.}}
\end{quote}

Ziel des Einschätzungsblockes ist das rasche Erkennen einer
vitalen Bedrohung, Überblick zu gewinnen und
Hauptbeschwerden/Leitsymptome zu ermitteln. Der
Einschätzungsblock muss in regelmäßigen Abständen sowie bei
Verdacht in angemessenem Umfang wiederholt werden, um eine
Verschlechterung des Patientenzustandes rechtzeitig zu
erkennen. 

\Text
{Zielsetzung und Struktur des Patientenmanagements nach \AASS}
{Ziel war es, ein möglichst schrittweise und geradlinig abarbeitenbares
System zu entwerfen, welches dennoch genügend Freiraum für spezielle
Situationen bietet und als \EE{Screening} für die wichtigsten vital bedrohlichen
Erkrankungen (und Verletzungen) geeignet ist. 
Das Modell folgt dem Grundsatz \Speech{\enquote{Treat
first what kills first}} und die traditionelle Reihenfolge
\enquote{Atemwege -- Atmung -- Kreislauf} wird natürlich
beibehalten. 
Der Einschätzungsblock ist eine Einheit und 
folgt einem durchgängigen Schema, bei dem die
Aktionen flüssig nacheinander abgearbeitet und nur durch
Sofortmaßnahmen unterbrochen werden.
\Commentary[Primary- und Secondary Survey]{Es wird formal
\EE{nicht} zwischen einer Primary- und Secondary Survey
unterschieden.
Wir haben festgestellt, dass im Dienstbetrieb die bisherigen
\E{Primary-/ Secondary-Modelle} eher nur bei schon von Anfang an als kritisch
erkennbaren Patienten (und hier vor allem bei Trauma-Patienten)
verwendet werden und die klassischen \enquote*{NACA-II-Patienten} nicht in den
Genuss eines systematisierten Screening gekommen sind. Daher haben wir
großen Wert darauf gelegt unser System so zu formulieren, dass es von
NACA I bis VII ohne große Anpassungen sinnvoll anwendbar ist
und haben auf die Unterscheidung der Primary- und Secondary
Survey bewusst verzichtet.}

\subparagraph{Aber was ist mit {\ldots}?}
Das ABCDE-Schema bzw. der Einschätzungsblock
\T{\Range{1}{E}} dient der Erkennung von vital bedrohten 
Patienten und dem Erkennen von Hauptbeschwerden und
Leitsymptomen, \E{es deckt nicht sämtliche
mögliche Untersuchungen ab!} Weiterführende Untersuchungen
und Anamneseschritte bzw. Maßnahmen werden \E{nach} dem
Einschätzungsblock nach Bedarf durchgeführt.
\Commentary[Was ist mit der General Impression?]{%
Traditionell werden einige Punkte vor dem Beginn des ABCDE-Blocks im
Rahmen einer \E{General Impression} zwar erwähnt, aber de
facto recht stiefmütterlich behandelt. Unser Konzept integriert die Punkte der
General Impression als \E{zum ABCDE gleichberechtigte Punkte} und betont
somit diese:
\EE{Szeneüberblick}, \EE{Eindruck}, \EE{Bewusstsein} und
\EE{Hauptbeschwerde} haben als Punkte \Range{\ABCDE{1}}{\ABCDE{4}} 
einen größeren Stellenwert als
bei vielen anderen Interpretationen des 
ABC(DE)-Schemas. Bei der Re-Evaluation werden diese Punkte 
ebenso erfasst.
Der Einschätzungsblock ist somit weitgehend kompatibel zum
simplen, alten \I{BAK-Schema} (Bewusstsein, Atmung, Kreislauf).
}

\subparagraph{Beispielhafte Maßnahmen}
Die angegebenen Maßnahmen stellen grundsätzlich nur
(häufige) Beispiele dar. Eine umfassende Aufzählung aller
möglichen Sofortmaßnahmen würde den Rahmen sprengen und wäre
sehr unübersichtlich. Außerdem führt eine taxative
Aufzählung zu einer Einschränkung der Flexibilität von
besonders qualifiziertem Personal. Weiters müssen Maßnahmen
immer \E{situationsgerecht} sein, und können aufgrund der
Vielzahl an Variablen nicht starr in Form eines Standard
Operating Procedure (SOP) festgeschrieben werden.
}
{}
{}
{}
{}




%%%%%%%%%%%%%%%%%%%%%%%%%%%%%%%%%%%%%%%%%%%%%%%%%%%%%%%%%%%
%%%%%%%%%%%%%%%%%%%%%%%%%%%%%%%%%%%%%%%%%%%%%%%%%%%%%%%%%%%
%%%%%%%%%%%%%%%%%%%%%%%%%%%%%%%%%%%%%%%%%%%%%%%%%%%%%%%%%%%
%%%%%%%%%%%%%%%%%%%%%%%%%%%%%%%%%%%%%%%%%%%%%%%%%%%%%%%%%%%
\clearpage
\subsection{{\FontKeyboardFamily{1}} Szeneüberblick und (Selbst-)Schutz}


\Dictum
[Mephistopheles. \\
In: Johann Wolfgang von Goethe: Faust. Der Tragödie erster Teil]
{Grau, teurer Freund, ist alle Theorie, \\
und Grün des Lebens goldner Baum.}

\TextSlide{\TxBeschreibung}{%
Kernpunkte sind die Fragen:
\begin{itemize}
  \item \Speech{Ist die Einsatzstelle sicher?}
  \item \Speech{Welche (Selbst-) Schutzmaßnhamen sind
  nötig?}
  \item \Speech{Sind weitere Ressourcen notwendig?}
  \item \Speech{Ist eine Lagemeldung erforderlich?}
\end{itemize}

Schon \E{vor} Eintreffen am Einsatzort beginnt der
Szeneüberblick. Das Einsatzpersonal muss sich die äußeren
Umstände des Einsatzes bewusst machen -- \EE{Wo} befindet
sich der Einsatzort -- in einer Wohnung in einem mehrstöckigen
Wohnhaus, in einem Einfamilienhaus oder auf einer
Verkehrsfläche? Wo befindet sich der \EE{Eingang}, wo gibt es
\EE{Parkmöglichkeiten}? Gibt es einen \EE{Aufzug} und wie
groß ist er? Können \EE{nachfolgende Kräfte} den Einsatzort
leicht finden? Wichtig ist das Erkennen von speziellen
Gefahren, wie z.\,B. Tiere (\Speech{\enquote{Vorsicht vor
dem Hund}}-Schilder, \E{Futternäpfe}), Hinweise auf
\E{Gasunfälle}, unsichere Gebäude(-teile), usw. Ggfs.
sind auch Fluchtwege einzuplanen.

Eventuell ist hier schon zu erkennen ob \EE{weitere
Ressourcen} benötigt werden, oder ob es sich um eine
\EE{Großschadenslage} oder einen \EE{Gefahrengutunfall}
handelt.
Ausserdem können Hinweise wahrgenommen
werden, ob ein \EE{Trauma} vorliegt. Bei Unfällen gibt die
Inspektion des Einsatzortes oft auch Aufschluss  über den
\EE{Unfallmechanismus} und -hergang. 
Wenn notwendig, muss der Leitstelle eine erste
\EE{Lagemeldung} gegeben werden.

}{%
\begin{itemize}
  \item Zufahrt, Gebäude 
  \item Gefahren
  \item Großschaden? Gefahrengut?
  \item Umstände, Situation, Umgebung, Ort, Zeit
  \item Trauma? Unfallmechanismus?
  \item Ggfs. Lagemeldung
\end{itemize}
}
{}
{}
{}

\begin{PictureSeriesWide}{}{Bilderserie ABCDE}
\PictureSeriesRowWide{3}
{./images/motal-aass/00800/helm1.jpg}
{Situation \SourcePicture{Motal}}
{./images/khd/gef-wind-rauch.jpg}
{Absperrung (hier deutlich verbesserungswürdig \ldots) \SourcePicture{Gabriel}}
{./images/KHD/warntafel-un-60-1710}
{Spezifische Gefahren erkennen}
{empty}
{}
\end{PictureSeriesWide}

\Text{Typische Maßnahmen}
{%
\begin{itemize}
  \item Absichern der Einsatzstelle
  \item Maßnahmen bei Gefahrgutunfällen (GAS-Maßnahmen, \FullPrettyRef{topic:gefahrengutunfall})
  \item Maßnahmen im Großschadensfall
  \item Lagemeldung
  \item Anlegen von Schutzausrüstung
  \item Nachforderung weiterer Kräfte
  \item Rettung aus einer Gefahrenzone
\end{itemize}
}

%%%%%%%%%%%%%%%%%%%%%%%%%%%%%%%%%%%%%%%%%%%%%%%%%%%%%%%%%%%
%%%%%%%%%%%%%%%%%%%%%%%%%%%%%%%%%%%%%%%%%%%%%%%%%%%%%%%%%%%
%%%%%%%%%%%%%%%%%%%%%%%%%%%%%%%%%%%%%%%%%%%%%%%%%%%%%%%%%%%
%%%%%%%%%%%%%%%%%%%%%%%%%%%%%%%%%%%%%%%%%%%%%%%%%%%%%%%%%%%
\clearpage
\subsection{{\FontKeyboardFamily{2}} Eindruck +
HWS-Stabilisierung (General Impression + Cervical Spine)} 

\TextSlide{{\TxBeschreibung}}{%
Hier zählt der Eindruck, den der Patient auf den Helfer
macht. Es werden v.\,a. beurteilt: 
\begin{compactitem}
  \item Alter, Geschlecht 
  \item Hautkolorit (Hautfärbung) 
  \item Haltung, Gesichtsausdruck, spontane Bewegungen 
  \item Sprache 
  \item Offensichtliche Verletzungen, Blutungen
  \item Sonstige Aufffälligkeiten
\end{compactitem}
Liegt der Verdacht auf eine Wirbelsäulenverletzung vor, wird
bereits hier mit der manuellen Stabilisierung der
Halswirbelsäule begonnen (\FullPrettyRef{m:einschaetzung-indikation-ws-immobilisation}). \index{HWS-Stabilisierung!Eindruck (ABCDE)}

}{%
\begin{itemize}
  \item Hautfarbe
  \item Gesichtsausdruck
  \item Haltung
  \item Sprache
  \item sonstige Auffälligkeiten
\end{itemize}
}
{}
{}
{}

\Text
{{\eye} Befunde}
{Auffallen kann eine \E{blasse  Haut}, \E{Zyanose},
\E{Schweiß}, \E{Schonhaltung},
\E{Atemnot} (Gestik, Aufstützen der Arme um die \E{Atemhilfsmuskulatur}
einzusetzen (\prettyref{topic:atemhilfsmuskulatur}),
Nasenflügeln, Einziehungen an den Rippen, \ldots{}), eine
verwaschene oder wirre \E{Sprache}, \E{Des\-orientiertheit},
\E{Krämpfe} u.\,v.\,a.\,m. Ebenso
können eventuell Hinweise auf ein \E{Trauma}, sichtbare
Verletzungen, starke Blutungen oder Hinweise auf eine
Wirbelsäulenverletzung erkannt werden.}
{\begin{compactitem}
    \item Blässe, Zyanose, Schweiß
    \item Schonhaltung, Atemnot
    \item Sprache, Desorientiertheit,
    Krämpfe
  \item Trauma, Verletzungen, Wirbelsäule, Blutungen
  \end{compactitem}}
{}
{}
{}

\begin{PictureSeriesWide}{}{Bilderserie ABCDE}
\PictureSeriesRowWide{3}
{./images/teaser/imgp0722_00500.jpg}
{Schädel-Hirn-Trauma: manuelle HWS-Immobilisierung  \SourcePicture{Gabriel}}
{./images/gabriel-sebastian-aass/UE2011FORTUNA-00651-0154pt-0300dpi.jpg}
{Zyanotische Lippen können schon frühzeitig auffallen \SourcePicture{Gabriel}}
{empty}
{}
{empty}
{}
\end{PictureSeriesWide}

\Text{Typische Maßnahmen}
{%
\begin{itemize}
  \item Manuelle HWS-Immobilisation
  \item Bewegungsverbot
  \item Stillen von starken Blutungen
\end{itemize}
}

%%%%%%%%%%%%%%%%%%%%%%%%%%%%%%%%%%%%%%%%%%%%%%%%%%%%%%%%%%%
%%%%%%%%%%%%%%%%%%%%%%%%%%%%%%%%%%%%%%%%%%%%%%%%%%%%%%%%%%%
%%%%%%%%%%%%%%%%%%%%%%%%%%%%%%%%%%%%%%%%%%%%%%%%%%%%%%%%%%%
%%%%%%%%%%%%%%%%%%%%%%%%%%%%%%%%%%%%%%%%%%%%%%%%%%%%%%%%%%%
\clearpage
\subsection{{\FontKeyboardFamily{3}} Bewusstsein
(Alertness)}



\label{topic:bewusstsein}
\index{Bewusstsein}



\TextSlide{\TxBeschreibung}
{\DefinitionInPara{Bewusstsein ist ein Oberbegriff für u.\,a. Wachheit,
Orientierung, Aufmerksamkeit, Auffassungsgabe, Denkverlauf und
Merkfähigkeit \cite{Pschy259}.}\ZFootnote{Die \enquote*{wahre} 
Definition von Bewusstsein ist seit
tausenden von Jahren ungeklärt. Auch an dieser Stelle wird
das Geheimnis nicht gelüftet werden, daher muss die obige
Definition für den Alltag ausreichen. Hinweisen
zufolge könnte die Antwort
eventuell \enquote{42} sein
\citep{UltimateHitchhikersGuide96hitchhikersguide}.}
Es ist eine grundlegende
\E{Vitalfunktion erster Ordnung} und ist besonders wichtig
hinsichtlich des Schutzes des Menschen vor Bedrohung. Ein bewusstseinsklarer
Mensch kann sich gegen Gefahren wehren, ein eingetrübter
oder bewusstloser Mensch kann dies schlecht oder gar nicht. 

Bei der Diagnostik beurteilt man 
den \EE{Bewusstseinsgrad} (\E{quantitativ};
\E{Wieviel} Bewusstsein ist vorhanden?) und
die \EE{Orientierung} (Qualität; \Speech{Wie gut funktioniert
das vorhandene Bewusstsein?} Zeitlich, örtlich, zur Situation und zur
Person). Die Orientierung wird unter
Punkt \ABCDE{D} beurteilt.
}{
\begin{itemize}
  \item Vitalfunktion 1. Ordnung
  \item Überlebenswichtig
  \item Beurteilung:
  \begin{itemize}
    \item Bewusstseinsgrad (quantitativ)
    \item Orientierung (qualitativ) $\rightarrow$ \ABCDE{D}
  \end{itemize}
\end{itemize}
}{}{}{}



%%%%%%%%%%%%%%%%%%%%%%%%%%%%%%%%%%%%%%%%%%%%%%%%%%%%%%%%%%%
\TextSlide{Bewusstseinsgrad}
{%
\index{Bewusstseinsgrad}
\label{topic:bewusstseinsstoerung-grade}
\index{Bewusstsein!Symptome}
%
Der Bewusstseinsgrad gibt Auskunft über die Quantität des
Bewusstseins (\Speech{\enquote{\EE{Wieviel} Bewusstsein
ist vorhanden?}}, vgl.
\prettyref{topic:bewusstseinsstoerung-grade}). 
Er muss bei jedem Patienten nach dem
\I{WASB-Schema} beurteilt
werden. 
%
\T{WASB} ist die Abkürzung für 
\EE{\underline{W}ach},
Reaktion auf \EE{\underline{A}nsprechen} oder
  leichte Weckreize,
Reaktion auf \EE{\underline{S}chmerzreiz} und
\EE{\underline{B}ewusstlos}. 
%
Dementsprechend kann ein Patient
\I{bewusstseinsklar}, \I{somnolent}, \I{soporös}
oder  \I{komatös} bzw. \I{bewusstlos} sein, vgl.
\prettyref{tab:bewusstseinsgrade}.
}{
\begin{itemize}
  \item W\,--\,A\,--\,S\,--\,B
  \begin{itemize}
    \item Wach \dotfill  \EE{klar}
    \item Ansprache \dotfill 
    \EE{somnolent}
    \item Schmerzreiz \dotfill  \EE{soporös}
    \item Bewusstlos \ldots  \EE{komatös}
  \end{itemize}
\end{itemize}
}{}{}{}


\Layout{57}
{}
{}
{}
{
\begin{tabularx}{\linewidth}{lllXl}
% \toprule
\addlinespace\toprule\addlinespace
\TabHeading{W}
&
\EE{W}ach
 &
\I{bewusstseinsklar}
 &
 wach
 &
\\\addlinespace
\TabHeading{A}
&
Reaktion auf \EE{A}nsprache
&
\noindent\I[Somnolenz]{somnolent}
&
schläfrig, aber leicht erweckbar.
 &
\\\addlinespace
\TabHeading{S}
&
Reaktion auf \EE{S}chmerzreiz
 &
\noindent\I[Sopor]{soporös}
 &
 kaum und nur mit erheblichen Aufwand (Schmerzreiz)
 erweckbar 
 &
\noindent\EE{Gefahr!}
\\\addlinespace
\TabHeading{B}
&
\EE{B}ewusstlosigkeit\index{Bewusstlosigkeit}
 &
\noindent\I[Koma]{komatös}
 &
 nicht erweckbar, ohne Bewusstsein
 &
\noindent\Ca{Lebensgefahr!}
\\\addlinespace\bottomrule
\end{tabularx}}
{Übersicht:
\EE{Bewusstseinsgrade}\label{tab:bewusstseinsgrade}.}
{}





\begin{table}[H]
\begin{WidePar}
\TableBaseModification
\Captionof{table}{GCS - Glasgow Coma Scale für den Erwachsenen.\label{tab:gcs}}

\begin{tabulary}{\linewidth}{LCCL}\addlinespace\toprule\addlinespace
\noindent ~ & \noindent\TabHeading{Max.} & \noindent\TabHeading{Punkte} &
\noindent\TabHeading{Reaktion}
\\\midrule
\noindent\TabSub{Augenöffnung} & 4 & 4 & spontan
\\
 & & 3 & auf Ansprache
\\
 & & 2 & Schmerzreiz
\\
 & & 1 & nicht
\\\midrule
\noindent\TabSub{Verbale Antwort} & 5 & 5 & klar
\\
 & & 4 & verwaschen
\\
 & & 3 & unzusammenhängende Worte
\\
 & & 2 & unartikuliert
\\
 & & 1 & nicht
\\\midrule
\noindent\TabSub{Motorische Reaktion} & 6 & 6 & spontan
\\
 & & 5 & gezielte Abwehrbewegungen
\\
 & & 4 & ungezielte Abwehrbewegungen
\\
 & & 3 & Beugung
\\
 & & 2 & Strecken
\\
 & & 1  & bewegt gar nicht
\\\bottomrule
\end{tabulary}
\end{WidePar}
\end{table}


%%%%%%%%%%%%%%%%%%%%%%%%%%%%%%%%%%%%%%%%%%%%%%%%%%%%%%%%%%%
\TextSlide{Bewusstseinsstörungen}
{%
\label{topic:bewusstseinsstoerung}%
\DefinitionInPara{Als Bewusstseinsstörungen bezeichnet man
insbesonders Störungen der
Wachheit und der Orientierung.}
Wie bereits oben erwähnt, ist das Bewusstsein wesentlich an
der \EE{Abwehr von Gefahren} beteiligt. Je mehr das
Bewusstsein gestört oder reduziert ist, desto gefährlicher ist das für
den Patienten. 

Ab einem gewissen Punkt können auch so wesentliche
Funktionen wie der \EE{Schluck-}, \EE{Würge-} bzw. \EE{Hustenreflex}
ausfallen, und es kann Mageninhalt, Speichel oder Blut in
die Lunge gelangen (\I{Aspiration}). Man spricht auch vom
\EE{Ausfall der
Schutzreflexe}\index{Schutzreflexe!Bewusstseinsstörung}.
Bei schweren Bewusstseinsstörungen kann es auch zum Erschlaffen
wichtiger Muskeln wie z.\,B. der Zunge kommen. Fällt
diese zurück, kann sie zurückfallen und den \EE{Atemweg
verlegen}.
Dabei ist manchmal ein dem \E{Schnarchen ähnliches}
Atemgeräusch zu vernehmen.
}{%
Gefahren:
\begin{itemize}
  \item Zurückfallen der Zunge
  \item Ausfall der Schutzreflexe
  \item Aspiration
\end{itemize}
}{}{}{}

\Fazit{Bewusstseinsstörungen können aufgrund des Ausfalls
der Schutzreflexe und des Zurückfallens der Zunge, sowie
der Unfähigkeit des Menschen, sich gegen Gefahren zu
wehren,
\EE{lebensgefährlich} sein!}




%%%%%%%%%%%%%%%%%%%%%%%%%%%%%%%%%%%%%%%%%%%%%%%%%%%%%%%%%%%
\TextSlide{Ursachen}
{%
Durch die Schädigung des zentralen Nervensystems kommt es zu
der Bewusstseinsstörung. Diese Schädigung kann direkt auf
das \Abk{ZNS} wirken (primäre Ursache) oder über Umwege
einen schädlichen Einfluss auf das \Abk{ZNS} haben (sekundäre
Ursache). \prettyref{tab:bewusstseinsstoerungen-ursachen}
listet typische Ursachen auf.

\begin{table}[H]
\TableBaseModification
\Captionof{table}{Typische Ursachen von
Bewusstseinsstörungen.\label{tab:bewusstseinsstoerungen-ursachen}}

\begin{tabularx}{\linewidth}[H]{XX}\addlinespace\toprule\addlinespace
\noindent\TabHeading{Primäre Ursachen} & \noindent\TabHeading{Sekundäre
Ursachen}
\\
{\tiny betreffen direkt das ZNS} & {\tiny betreffen über
Umwege das ZNS} 
\\\addlinespace\cmidrule(lr){1-1}\cmidrule(lr){2-2}%\addlinespace
\begin{itemize}
    \item Verletzungen des ZNS
    \item Blutungen, Hirnschwellung
    \item Schlaganfall
    \item Epileptischer Anfall
    \item Entzündungen
    \item \Z{Tumor}
    \item \ldots
\end{itemize}
 &
 \begin{itemize}
   \item Störungen der \ce{O2}-Versorgung:
   \begin{compactitem}
     \item Atemstörungen  (Hypoxie)
     \item Herz-Kreislaufstörungen
   \end{compactitem}
   \item Störungen des Wasserhaushaltes und
   Elektrolytentgleisungen
   \item Stoffwechselstörungen
   \begin{compactitem}
     \item Zuckerhaushalt
     \item Leberversagen
   \end{compactitem}
   \item Vergiftungen:
   \begin{compactitem}
     \item Alkohol
     \item Medikamente, Drogen
     \item Reizgase,  Lösungsmittel ...
   \end{compactitem}
   \item \ldots
 \end{itemize}
\\\bottomrule
\end{tabularx}
\end{table}
}
{
\begin{itemize}
  \item Direkt das \Abk{ZNS} betreffend (primär)
  \item Indirekt das \Abk{ZNS} betreffend (sekundär)
\end{itemize}

\prettyref{tab:bewusstseinsstoerungen-ursachen}.
}{}{}{}




% M %%%%%%%%%%%%%%%%%%%%%%%%%%%%%%%%%%%%%%%%%%%%%%%%%%% M %
% Bewusstseinseintrübung
\ProcedurePrintDefault{MR40001B}


% M %%%%%%%%%%%%%%%%%%%%%%%%%%%%%%%%%%%%%%%%%%%%%%%%%%% M %
% Bewusstlose und soporöse Patienten
\ProcedurePrintDefault{MR40021B}



% \begin{PictureSeriesWide}{}{Bilderserie ABCDE}
% \PictureSeriesRowWide{3}
% {./images/gabriel-sebastian-aass/Gallenblase-Palpation-001-00800.jpg}
% {ABCDE-Bild}
% {./images/gabriel-sebastian-aass/Gallenblase-Palpation-001-00800.jpg}
% {ABCDE-Bild}
% {./images/gabriel-sebastian-aass/Gallenblase-Palpation-001-00800.jpg}
% {ABCDE-Bild}
% {empty}
% {}
% \end{PictureSeriesWide}

\Text{Typische Maßnahmen}
{%
\begin{itemize}
  \item \TxMassVit 
  \begin{compactitem}
    \item Absaugbereitschaft!
  \end{compactitem}
  \item Stabile Seitenlage nach \ABCDE{B}
\end{itemize}
}

%%%%%%%%%%%%%%%%%%%%%%%%%%%%%%%%%%%%%%%%%%%%%%%%%%%%%%%%%%%
%%%%%%%%%%%%%%%%%%%%%%%%%%%%%%%%%%%%%%%%%%%%%%%%%%%%%%%%%%%
%%%%%%%%%%%%%%%%%%%%%%%%%%%%%%%%%%%%%%%%%%%%%%%%%%%%%%%%%%%
%%%%%%%%%%%%%%%%%%%%%%%%%%%%%%%%%%%%%%%%%%%%%%%%%%%%%%%%%%%
% \clearpage
\subsection{{\FontKeyboardFamily{4}} Hauptbeschwerde (Main
Complaint)} 

\TextSlide{What's up, man?}{%
Die \EE{Hauptbeschwerde} des Patienten, bzw. der \EE{Grund
der Berufung}, sowie Schmerzen und andere für den Patienten
dringliche Symptome müssen herausgefunden werden. I.\,d.\,R.
wird hier auf den \enquote{S}-Teil des \Abk{SAMPLE}
vorgegriffen.
Es bietet sich die Frage an:
\begin{quote}
\Speech{\enquote{Was können wir für sie tun?}}
\end{quote}
Eventuell ist auch die Umgebung (Angehörige, Dritte) einzubeziehen
(Fremdanamnese).
}{%
\begin{itemize}
  \item Eigene Wahrnehmung, Patientengespräch, Angehörige
\end{itemize}
\begin{compactitem}
  \item[\eye] Grund der Berufung, Schmerzen und andere für
  den Patienten dringliche Symptome, Trauma
\end{compactitem}
}


% \begin{PictureSeriesWide}{}{Bilderserie ABCDE}
% \PictureSeriesRowWide{3}
% {./images/gabriel-sebastian-aass/Gallenblase-Palpation-001-00800.jpg}
% {ABCDE-Bild}
% {./images/gabriel-sebastian-aass/Gallenblase-Palpation-001-00800.jpg}
% {ABCDE-Bild}
% {./images/gabriel-sebastian-aass/Gallenblase-Palpation-001-00800.jpg}
% {ABCDE-Bild}
% {empty}
% {}
% \end{PictureSeriesWide}

\Text{Typische Maßnahmen}
{%
\begin{itemize}
  \item Gezielte Kurzanamnese
\end{itemize}
}

%%%%%%%%%%%%%%%%%%%%%%%%%%%%%%%%%%%%%%%%%%%%%%%%%%%%%%%%%%%
%%%%%%%%%%%%%%%%%%%%%%%%%%%%%%%%%%%%%%%%%%%%%%%%%%%%%%%%%%%
%%%%%%%%%%%%%%%%%%%%%%%%%%%%%%%%%%%%%%%%%%%%%%%%%%%%%%%%%%%
%%%%%%%%%%%%%%%%%%%%%%%%%%%%%%%%%%%%%%%%%%%%%%%%%%%%%%%%%%%
\clearpage
\subsection{{\FontKeyboardFamily{A}} Atemweg (Airway)}

\label{topic:atemweg}
\label{topic:airway}


\TextSlide{Vom Gurgeln und Schnarchen}{%
Unter Punkt
A wird eingeschätzt, ob der Atemweg {frei},
{verlegt} oder {gefährdet} ist.
Dazu werden zuerst die \E{Mundhöhle}, die \E{Nase} und der
\E{Hals} \EE{inspiziert}. Liegt ein Verdacht auf eine
Verletzung vor, werden Nase und Hals auch \E{abgetastet}.
Weiters werden die \EE{Atemgeräusche} (ohne Hilfsmittel) beurteilt. 

}{%
\begin{itemize}
  \item Inspektion \& ggfs. Abtasten von \EE{Mundhöhle},
  \EE{Nase} und \EE{Hals}
\end{itemize}
}
{}
{}
{}

\TextSlide
{{\eye} Befunde}
{Speziell \E{pfeifende}, \E{gurgelnde} oder \E{schnarchende}
Atemgeräusche können ein Hinweis auf eine Verlegung sein.
Oft gibt auch die \EE{Gestik} und Haltung des Patienten
Hinweise auf eine Verlegung. 
Ein \I[Atemweg!gefährdeter]{gefährdeter Atemweg} liegt vor, wenn dieser noch nicht
verlegt, jedoch aufgrund des Patientenzustandes, der
Diagnose oder anderer Gegebenheiten eine baldige Verlegung
wahrscheinlich ist (Verbrennungen des Brustkorbes, allergische Reaktion, \ldots).

Der Atemweg kann \EE{frei},
\EE{verlegt} oder \EE{gefährdet} sein.
}
{
\begin{itemize}
  \item Frei, verlegt, gefährdet
  \item Atemgeräusche: Pfeifen, Gurgeln, Schnarchen;
  Gestik
\end{itemize}
}
{}
{}
{}

% \begin{PictureSeriesWide}{}{Bilderserie ABCDE}
% \PictureSeriesRowWide{3}
% {./images/ASBOE-Grassl-gjh-1/IMG_0288-00441pt}
% {Atemkonmtrolle bei überstrecktem Kopf: Sehen, hören, fühlen! \SourcePicture{ASBÖ, Jürgen Grassl, jgh / MfG}}
% {./images/ASBOE-Grassl-gjh-1/IMG_0268-00441pt}
% {Überstrecken des Kopfes \SourcePicture{ASBÖ, Jürgen Grassl, jgh / MfG}}
% {./images/ASBOE-Grassl-gjh-1/IMG_0351-copy2-00441pt}
% {Esmarch-Handgriff \SourcePicture{ASBÖ, Jürgen Grassl, jgh / MfG}}
% {empty}
% {}
% \PictureSeriesRowWide{3}
% {./images/hirtler-aass/Fremdkoerper.pdf}
% {Fremdkörper in Luft- und Speiseröhre \SourcePicture{Hirtler}}
% {./images/ASBOE-Grassl-gjh-1/verschlucken_03-00441pt}
% {Fremdkörper verschluckt}
% {./images/ASBOE-Grassl-gjh-1/verschlucken_08-00441pt}
% {Schulterklopfen}
% {empty}
% {}
% \end{PictureSeriesWide}

% TODO Grassl Foto: A, Fremdkörper

\Text{Typische Maßnahmen}
{%
\begin{itemize}
  \item Überstrecken des Kopfes
  \item Esmarch-Handgriff
  \item Stabile Seitenlage nach \ABCDE{B}
  \item Absaugung
  \item Spezielle Maßnahmen bei Atemwegsverlegung (\FullPrettyRef{topic:mechanische-atemwegsverlegung})
  \item \TxMassVit 
  \begin{compactitem}
    \item Absaugbereitschaft!
  \end{compactitem}
\end{itemize}
}

%%%%%%%%%%%%%%%%%%%%%%%%%%%%%%%%%%%%%%%%%%%%%%%%%%%%%%%%%%%
%%%%%%%%%%%%%%%%%%%%%%%%%%%%%%%%%%%%%%%%%%%%%%%%%%%%%%%%%%%
%%%%%%%%%%%%%%%%%%%%%%%%%%%%%%%%%%%%%%%%%%%%%%%%%%%%%%%%%%%
%%%%%%%%%%%%%%%%%%%%%%%%%%%%%%%%%%%%%%%%%%%%%%%%%%%%%%%%%%%
\clearpage
\subsection{{\FontKeyboardFamily{B}} Atmung (Breathing)}

\label{topic:atmung}
\label{topic:breathing}

\TextSlide{Take a deep breath \ldots}{%
\index{Atemprobe}
\index{Atmung!Einschätzungsblock}
Als nächstes wird die Funktion der \E{Atmung} beurteilt.
Dazu wird die Frequenz  und Tiefe der Atmung
eingeschätzt. Dabei wird der Thorax betrachtet und die
Atmung beobachtet, während der Einatmung wird auf das
\E{gleichseitige} Heben des Brustkorbes geachtet. 
Während der Atmung wird auf \EE{Atemgeräusche}
geachtet. 

Bestimmte charakteristische Zeichen können auf
Atmungsprobleme hinweisen. Dazu gehören z.\,B.
die \E{Hautfarbe},
pathologische \E{Bauch- und Thoraxbewegungen},
\E{Anstrengungen} beim Atmen, 
Einsatz der \E{Atemhilfsmuskulatur},
\E{Einziehungen} an den Rippen, 
\E{Nasenflügeln} (bei Kindern), sowie
andere Anzeichen der Atemnot. 

\Z{\E{Speziell ausgebildetes} Personal kann mittels eines
Stethoskops die oberen und unteren Atemwege abhören
(auskultieren) und damit Störungen des Atemweges und der
Atmung besser einschätzen. 
Ebenso ist diesem bei Verdacht auf das Vorliegen einer 
mechanisch bedingten Atemstörung das Vorziehen des Abtastens 
des Thorax möglich.}

}{%
\begin{itemize}
  \item Schätzen: \E{Frequenz} und \E{Tiefe}
  \item Atemgeräusche
  \item Inspektion von Thorax inkl. Atemprobe
  \item Anzeichen: \\
  Hautfarbe, Bauch-/Thorax\-bewegungen,
  Anstrengung beim Atmen, Atemhilfsmuskulatur, Zeichen der
  Atemnot
  %   (\prettyref{topic:atemnot-zeichen-kinder})
  \item Heimsauerstoff
  \item \Z{(Auskultation: Seitenvergleich, Lungenbasen)}
\end{itemize}
}
{}
{}
{}



\subsubsection{Physiologie der Atmung}

% STEU 2009-06-28: Trennung Oxygenation und Ventilation
\A{STEU 2009-06-28: Trennung Oxygenation und Ventilation}

\TextSlide{Atemfrequenz, -tiefe und -volumen}
{%
\index{Atemfrequenz}
\index{AF}
\index{Atemzugsvolumen}
\index{AZV}
\index{Atemminutenvolumen}  
\index{AMV} 
\label{topic:atemvolumina}
%
Es gibt drei elementare Kennzahlen, welche die wichtigsten Atmungsparameter
beschreiben
\cite{KlinkePhysio3,SchmidtPhysio28,DRAn3,StriebelAn7}:

\begin{enumerate}
  \item Die \T{Atemfrequenz} (\T{AF}) gibt die Anzahl der
  Atemzüge pro Minute an. Die normale Atemfrequenz  beträgt beim Erwachsenen
  in Ruhe \VonBis{12}{16}\PerMin*.
  \item Das \T{Atemzugsvolumen} (\T{AZV}) gibt die Menge
  Luft an, die \E{pro Atemzug} eingeatmet wird. Das normale
  Atemzugsvolumen  eines Erwachsenen beträgt ca.
  500\Ml*.\ZFootnote{Die Einstellungen eines Beatmungsgeräts
  können von den Normalwerten deutlich abweichen, da man
  z.\,B. manchmal
  vorsichtiger, bzw. lungenschonender, beatmen möchte.}
  \item Das \T{Atemminutenvolumen} (\T{AMV}) gibt die Menge
  an \E{pro Minute} eingeatmeter Luft an. 
  Es kann aus der
  Atemfrequenz und dem Atemzugsvolumen \E{errechnet} werden:
  {\FontSansNarrowFamily
  \begin{equation}
    \mbox{AMV} = \mbox{AZV} \times \mbox{AF}
  \end{equation}}
\end{enumerate}
Die oben angegebenen Normalwerte ergeben also ein normales
  Atemminutenvolumen von  \VonBis{6}{8}\LPerMin*. 

}{
\begin{labeling}[:]{AMV}
    \item[AF] Atemfrequenz,
    \EE{\VonBis{12}{16}\PerMin*}
    \item[AZV] Atemzugsvolumen,
    \EE{500\Ml*}
    \item[AMV] Atemminutenvolumen,
    \EE{\VonBis{6}{8}\LPerMin*}
\end{labeling}

\begin{center}
{\FontSansNarrowFamily
$\mbox{AMV} = \mbox{AZV} \times \mbox{AF}$
}
\end{center}
}{}{}{}


\begin{table}[H]
\Captionof{table}{Die genauen Normalwerte sind abhängig von den Faktoren
Alter, Geschlecht und Größe: Atemfrequenz und -zugvolumen im
Altersvergleich (Richtwerte)}

\begin{tabular}{lrr}\addlinespace\toprule\addlinespace
 & \noindent\TabHeading{AF} & \noindent\TabHeading{AZV}
\\\addlinespace\midrule\addlinespace
\TabSub{Neugeborenes} & 40\PerMin* & \VonBis{20}{40}\Ml*
\\\addlinespace 
\TabSub{Kleinkind} & 25\PerMin* & \VonBis{100}{200}\Ml* 
\\\addlinespace
\TabSub{Schulkind} & 20\PerMin* & \VonBis{200}{400}\Ml* 
\\\addlinespace
\TabSub{Erwachsener} & \VonBis{12}{16}\PerMin*  & ca. 500\Ml*
\\\addlinespace\bottomrule
\end{tabular}
\end{table}

\Fazit{AMV = AZV $\ast$ AF, beim Erwachsenen  ca.
\VonBis{6}{8}\Liter*}

Für
eine ausreichende Atmung ist auch eine ausreichende
\ce{O2}-Konzentration der Umgebungsluft notwendig, denn ohne
Sauerstoff in der Einatemluft ist die beste Atmung nutzlos.
Der normale \ce{O2}-Gehalt der Umgebungsluft
beträgt 21\PC*.



\TextSlide{Totraum}
{%
\index{Totraum}%
Unter dem Totraum versteht man den Teil der Atemwege der
nicht am Gasaustausch teilnimmt. Dies ist im wesentlichen
der gesamte Weg von der Nasenspitze bis kurz vor den
Lungenbläschen (Alveolen). Die Luft muss diesen Totraum
passieren um in die Lungenbläschen zu gelangen! Der Totraum
beträgt beim Erwachsenen ca. 150\Ml*. \cite{KlinkePhysio3}

% \index{Totraumventilation}
Bei der \T[Totraum!-ventilation]{Totraumventilation} ist das
Atemzugsvolumen so klein, dass die Frischluft im
Totraum verbleibt und nicht in die Lungenbläschen gelangt: Es gelangt 
folglich kein Sauerstoff ins Blut.
\label{topic:schnappatmung}
Bei der \T{Schnappatmung} kann
genau dies passieren: Auf den ersten Blick sieht es so aus
als würde der Patient \enquote*{schnappend} atmen, in
Wirklichkeit ist das Atemzugsvolumen so gering dass es zu einer Totraumventilation
kommt: Die Atmung ist nicht ausreichend.
}{
\begin{itemize}
    \item Totraum: Teil der Atemwege, der nicht am
    Gassaustausch teilnimmt
    \item Totraumventilation
    \item Schnappatmung = nicht ausreichend
\end{itemize}
}{}{}{}



\Cave{Schnappatmung = Keine suffiziente Atmung}




\TextSlide{Steuerung der Atmung}
{%
Die Atmung wird durch das \EE{Atemzentrum im Hirnstamm}
gesteuert. Daneben gibt es Sensoren im Körper, die den
Gasgehalt des Blutes beobachten.
Beim gesunden Patienten ist der \EE{Atemantrieb} vom
\EE{\ce{CO2}-Gehalt} des Blutes abhängig. Der
Sauerstoffanteil spielt nur eine untergeordnete Rolle! Unter
\FullPrettyRef{topic:saeure-basen-haushalt}, wird
ausgeführt, dass das \ce{CO2} eine wichtige Rolle im
Säure-Basen-Haushalt des Körpers einnimmt.

Wenn bei \EE{bestimmten Krankheitsbildern} die
\E{\ce{CO2}-Konzentration ständig hoch} ist, \E{gewöhnt}
sich der Körper an diesen Zustand. Er beginnt dann den
\E{Atemantrieb direkt über den \ce{O2}}-Gehalt im Blut zu
regeln.
Dies kann \EE{Probleme} machen: Wird bei so einem Patienten
der \ce{O2}-Gehalt plötzlich künstlich erhöht, fehlt der
Atemantrieb und der Patient kann \EE{Bewusstseinsstörungen}
(\I[\ce{CO2}-Narkose@CO2-Narkose]{\ce{CO2}-Narkose}) und
einen \EE{Atemstillstand} erleiden.
Das wichtigste Krankheitsbild bei dem so etwas vorkommen
kann, ist die \Abk{COPD} (\FullPrettyRef{topic:copd}).
}{
\begin{itemize}
  \item Atemzentrum im Hirnstamm
  \item Atemantrieb normalerweise durch \EE{\ce{CO2}}
  \item Krankhaft:
  \begin{itemize}
    \item Wenn \ce{CO2}-Konzentration dauerhaft erhöht
    \Sign{→} Regelung durch
    \ce{O2}-Gehalt.
    \item Problem bei Sauerstoffgabe
    \item V.\,a. bei \Abk{COPD} (\ref{topic:copd})
  \end{itemize}
\end{itemize}
}{}{}{}

\Fazit{Normalerweise wird der Atemantrieb über den
\ce{CO2}-Gehalt des Blutes gesteuert.}




%%%%%%%%%%%%%%%%%%%%%%%%%%%%%%%%%%%%%%%%%%%%%%%%%%%%%%%%%%%
%%%%%%%%%%%%%%%%%%%%%%%%%%%%%%%%%%%%%%%%%%%%%%%%%%%%%%%%%%%
%%%%%%%%%%%%%%%%%%%%%%%%%%%%%%%%%%%%%%%%%%%%%%%%%%%%%%%%%%%
\subsubsection{Störungen der Atmung}
\label{topic:atemstörungen}
\label{topic:atemstillstand}

\TextSlide{Einleitung}
{%
\index{Ateminsuffizienz}%
\index{Insuffizienz!Atem-}%
\index{Dyspnoe}%
\index{Apnoe}%
\index{Bradypnoe}%
\index{Tachypnoe}%
Störungen der Atmung können aus vielen Ursachen eintreten.
\FullPrettyRef{tab:atemstoerungen-ursachen}, beschreibt die
wichtigsten Pathomechanismen und Störungen.
Der Begriff \T{Ateminsuffizienz} beschreibt allgemein eine
\EE{nicht ausreichende Atmung}, unabhängig von der
Ursache der Störung. Dies führt zu einem \ce{O2}-Mangel
(Hypoxie). Oft hat der Patient das Gefühl nicht genug Luft
zu bekommen und leidet an \EE{Atemnot} (\T{Dyspnoe}).
Der \T{Atemstillstand} (\T{Apnoe}) ist die extremste Form
der Ateminsuffizenz. Störungen der Atemgeschwindigkeit
werden als \I{Bradypnoe} (zu langsam) und als \I{Tachypnoe}
(zu schnell) bezeichnet.
}{
\begin{itemize}
  \item Ateminsuffizienz \ldots{} nicht ausreichende Atmung
  \item Hypoxie \ldots{} \ce{O2}-Mangel
  \item Atemnot (Dyspnoe)
  \item Atemstillstand (Apnoe)
  \item Bradypnoe (zu langsam)
  \item Tachypnoe (zu schnell)
\end{itemize}
}{}{}{}


\TextSlide{Gefahr der Ateminsuffizienz}
{%
Bei einer Ateminsuffizienz sinkt der \EE{\ce{O2}-Gehalt} im
Blut. Es kommt zu einer \EE{Hypoxie} und in weiterer Folge
zu einer Unterversorgung lebenswichtiger Organe wie Herz und
Hirn.
Doch auch die \EE{\ce{CO2}-Abgabe} ist gestört und es kommt
zu einer Anhäufung von \ce{CO2} im Blut. In weiterer Folge
kann es zu einer \E{atmungsbedingten Übersäuerung}
(\EE{Respiratorische Azidose},
\FullPrettyRef{topic:azidose-respiratorische}) kommen.
}{
\begin{itemize}
    \item \ce{O2}-Gehalt im Blut ↓ (Hypoxie)
    \item \ce{CO2}-Gehalt im Blut ↑
    \item Respiratorische Azidose
    (\FullPrettyRef{topic:azidose-respiratorische})
\end{itemize}
}{}{}{}

\AuthNotes{GABG: pathologische Atmungstypen inkl.
Zeichnung einbinden? Kussmaul, etc.}



\TextSlide{Ursachen}
{%
\prettyref{tab:atemstoerungen-ursachen}.
}{\prettyref{tab:atemstoerungen-ursachen}.}{}{}{}

\Layout{57}
{}
{}
{}
{\begin{tabulary}{\linewidth}[H]{LL}\addlinespace\toprule\addlinespace
\noindent\TabHeading{Störung des/der}  &
\noindent\TabHeading{Vorkommen und Erkrankungen}
\\\addlinespace\midrule\addlinespace
\noindent\TabSub{\ce{O2}-Angebotes}
 &
Gärkeller (\ce{CO2}\,${\uparrow}$, \ce{O2}\,${\downarrow}$)

\enquote{dünne Luft} (z.\,B. Hochgebirge;
niedriger Sauerstoffpartialdruck)

Defekte Gastherme (\ce{CO}\,${\uparrow}$,
\ce{O2}\,${\downarrow}$)

Ertrinken
\\\addlinespace
\noindent\TabSub{\ce{O2}-Diffusion}
 &
Lungenödem

Lungenentzündung

Lungenembolie
\\\addlinespace
\noindent\TabSub{Atemmechanik}
 &
Verlegung der oberen Atemwege (Zunge, Erbrochenes, Laryngospasmus,
         Glottisödem, Bolus)

Verlegung der unteren Atemwege (Bronchitis, Asthma,
Lungenödem)

Verminderung der Dehnbarkeit der Thoraxwand oder des Lungengewebes
     (Rippenfraktur, Pneumothorax, Pleuraerguß, Emphysem)
\\\addlinespace
\noindent\TabSub{Neuro\-muskulären Atemregulation}
 &
SHT (Schädel Hirn Trauma)

Schlaganfall (Insult)

Vergiftungen

Tumore

Hohe Querschnittslähmung (Rückenmark durchtrennt)

Muskelerkrankungen

Nervenerkrankung

Stoffwechselstörungen
\\\addlinespace
\noindent\TabSub{Sonstige Regulationsstörung}
 &
 Psychogen (z.\,B. Panikattacke)
\\\addlinespace\bottomrule
\end{tabulary}}
{Ursachen für
Atemstörungen.\label{tab:atemstoerungen-ursachen}}
{}


\TextSlide
{{\eye} Befunde}
{Atemstörungen können viele verschiedene Symptome
verursachen. \FullPrettyRef{tab:atmung-symptome}, listet
verschiedene Symptome auf. Neben technischen Geräten ist vor
allem \E{sehen},
\E{hören}, \E{fühlen} und ein \E{aufmerksamer Blick}
gefordert.
Die Abgrenzung der Atemstörung zur Kreislaufstörung ist oft
nicht einfach. Eine Störung des Kreislaufes hat meist auch
einen verminderten Sauerstoff\E{transport} zur Folge, daher
kommt es dabei häufig {auch} zur Atemnot.

Häufig erkennbar sind z.\,B. eine \EE{Zyanose}, 
Klage über subjektive \EE{Atemnot}, 
eine zu schnelle
oder langsame Atmung (Tachy-/Bradypnoe; 
\EE{AF <\,8\PerMin* oder >\,30\PerMin*}), 
eine auffällige \EE{Atemtiefe} (flach, tief, evtl. \EE{Schnappatmung}), 
pathologische Atemmuster, 
eine paradoxe oder überhaupt fehlende Atmung oder 
andere Symptome der Atemnot
(\prettyref{tab:atmung-symptome}).
\Z{Bei der Auskultation kann eine ungleichseitige Belüftung der 
Lungen festgestellt werden, 
sowie diverse Atemnebengeräusche 
(feuchte Rasselgeräusche, spastische Atemgeräusche).}

Die \EE{Sauerstoffsättigung} kann vermindert 
oder hoch sein (Kohlenmonoxid, Hyperventilationssyndrom). 
Ebenso wichtig sind Umstände wie das Vorhandensein 
von Heimsauerstoff
oder anderen Atemhilfen.
}
{\begin{itemize}
  \item Zyanose, Atemnot
  \item Tachy-/Bradypnoe (AF <\,8\PerMin* oder >\,30\PerMin*)
  \item Atemtiefe (flach, tief, Schnappatmung)
  \item Path. Atemmuster
  \item Paradoxe, fehlende Atmung
  \item Sauerstoffsättigung
  \item \Z{Auskultation: Belüftung, feuchte
  Rasselgaräusche, Spastik}
  \item Heimsauerstoff, Atemhilfen
  \item \FullPrettyRef{tab:atmung-symptome}
\end{itemize}
}
{}
{}
{}


\Layout{57}
{}
{}
{}
{\begin{tabulary}{\linewidth}[H]{LLLL}
\toprule\addlinespace
\noindent\TabHeading{Kriterium}
 &
\noindent\TabHeading{Befund}
 &
\noindent\TabHeading{Beschreibung}
 &
%\noindent\TabHeading{Anm.}
\\\addlinespace\midrule\addlinespace
\noindent\TabSub{Beschwerde}
 &
Atemnot (Dyspnoe) \index{Dyspnoe}
 &
Leitsymptom
 &
\\\addlinespace
\noindent\TabSub{Atemgeräusch}
 &
Brodeln \index{Atemgeräusch!brodelndes}
 &
Blubbern, klassisch für Lungenödem
\index{Atemgeräusch!blubberndes} &
\\\addlinespace
 &
Stridor \index{Stridor}
 &
 &
\\\addlinespace
 &
Brummen, Giemen
\index{Atemgeräusch!brummendes}\index{Atemgeräusch!giemendes}
& &
\\\addlinespace
 &
Rasselgeräusche \index{Rasselgeräusche}
 &
 &
\\\addlinespace
\noindent\TabSub{Frequenz}
 &
Beschleunigt
 &
Tachypnoe \index{Tachypnoe}
 &
\\\addlinespace
 &
Verlangsamt
 &
Bradypnoe \index{Bradypnoe}
 &
\\\addlinespace
\noindent\TabSub{Atemzugs\-volumen}
 &
Schnappatmung
 &
\FullPrettyRef{topic:schnappatmung}
 &
\\\addlinespace
 &
Flache Atmung
 &
 &
\\\addlinespace
 &
Tiefe Atmung
 &
Z.\,B. Kußmaul'sche Atmung
 &
\\\addlinespace
\noindent\TabSub{Hautfarbe}
 &
Blass
 &
Normal, Anämie, Blutverlust
 &
\\\addlinespace
 &
Rosig
 &
Normal, \ce{CO}-Vergiftung
 &
\\\addlinespace
 &
Bläulich 
 &
Zyanose: Hypoxie!
 &
\\\addlinespace
\noindent\TabSub{Körperlich}
 &
Einziehungen an den Rippen
&
Einsatz der Atemhilfsmuskulatur
&
\\\addlinespace
&
Aufstützen
&
Einsatz der Atemhilfsmuskulatur
&
\\\addlinespace
&
Nasenflügeln \index{Nasenflügeln}
&
Atemnot, besonders bei Kleinkindern
&
\\\addlinespace
&
Aufrechte Position
&
&
\\\addlinespace
&
\T[Atmung!paradoxe]{Paradoxe Atmung}
&
Brustkorb \E{senkt} sich bei Einatmung, z.\,B. bei
Serienrippenfraktur 
&
\\\addlinespace
\noindent\TabSub{Geräte}
&
Pulsoxymetrie
&
\FullPrettyRef{topic:pulsoxymetrie}
&
\\\addlinespace
&
\Z{Kapnometrie}
&
\Z{\ce{CO2}-Messung}
\\\addlinespace\bottomrule
\end{tabulary}}
{Symptome: Atmung.\label{tab:atmung-symptome}}
{}





% M %%%%%%%%%%%%%%%%%%%%%%%%%%%%%%%%%%%%%%%%%%%%%%%%%%% M %
% Insuffiziente Atmung
\ProcedurePrintDefault{MJ96091C}


% M %%%%%%%%%%%%%%%%%%%%%%%%%%%%%%%%%%%%%%%%%%%%%%%%%%% M %
% Atemstillstand
\ProcedurePrintDefault{MR09020B}





\begin{PictureSeriesWide}{}{Bilderserie ABCDE}
\PictureSeriesRowWide{3}
{./images/hirtler-aass/Atmung.pdf}
{Hirtler}
{./images/gabriel-sebastian-aass/UE2011FORTUNA-00651-0154pt-0300dpi.jpg}
{Zyanotische Lippen \SourcePicture{Gabriel}}
% {./images/ASBOE-Grassl-gjh-1/auscultation-00441pt}
% {ABCDE-Bild}
{empty}
{}
{empty}
{}
\end{PictureSeriesWide}

% TODO Grassl Foto: Auskultation

\Text{Typische Maßnahmen}
{%
\begin{itemize}
  \item Sauerstoffgabe
  \item Lippenbremse
  \item Beengende Kleidung öffnen
  \item Lagerung mit erhöhtem Oberkörper
  \item Atemkommandos
  \item Beatmung
  \item \TxMassVit
\end{itemize}
}


%%%%%%%%%%%%%%%%%%%%%%%%%%%%%%%%%%%%%%%%%%%%%%%%%%%%%%%%%%%
%%%%%%%%%%%%%%%%%%%%%%%%%%%%%%%%%%%%%%%%%%%%%%%%%%%%%%%%%%%
%%%%%%%%%%%%%%%%%%%%%%%%%%%%%%%%%%%%%%%%%%%%%%%%%%%%%%%%%%%
%%%%%%%%%%%%%%%%%%%%%%%%%%%%%%%%%%%%%%%%%%%%%%%%%%%%%%%%%%%
\clearpage
\subsection{{\FontKeyboardFamily{C}} Kreislauf
(Circulation) und Schnelle Trauma-Untersuchung (STU)}

\label{topic:kreislauf}
\label{topic:circulation}


% \begin{PictureSeriesWide}{}{Bilderserie ABCDE}
% \PictureSeriesRowWide{3}
% {./images/hirtler-aass/Kreislauf-Schema.pdf}
% {Hirtler}
% {./images/gabriel-sebastian-aass/Gallenblase-Palpation-001-00800.jpg}
% {ABCDE-Bild}
% {./images/gabriel-sebastian-aass/Gallenblase-Palpation-001-00800.jpg}
% {ABCDE-Bild}
% {empty}
% {}
% \end{PictureSeriesWide}




\TextSlide
{Blut als Sauerstofftransporteur}
{Durch das Kreislaufsystem wird sauerstoff- und
nährstoffreiches Blut gepumpt. Ausschlaggebend für die
Sauerstoffversorgung ist der \I{Sauerstoffgehalt} des
Blutes. Dieser wird wesentlich durch den Gasaustausch, aber
auch durch die Aufnahmekapazität des Blutes bestimmt.
Wieviel Sauerstoff das Blut aufnehmen kann, hängt vor allem
von der Anzahl der roten Blutkörperchen (\I{Erythrozyten})
bzw. der Konzentration des Blutfarbstoffes \I{Hämoglobin}
ab.}
{
\begin{itemize}
  \item Rote Blutkörpechen (Erythrozyten) und Hämoglobin
\end{itemize}
}
{}
{}
{}



\Z{Im arteriellen Blut sind pro Liter ca 0,2\Liter*
Sauerstoff enthalten. Im Normalfall werden davon nur ca. 25\PC* vom
Körper aufgenommen, der Rest gelangt ungenutzt in das venöse
System zurück.
\cite{Silbernagl3-Kreislaufsystem}}


\TextSlide
{Herzleistung und Blutdruck}
{Das Herz schlägt beim Erwachsenen mit einer
\I{Herzfrequenz} (\I{HF}, bzw. \I{Pulsfrequenz}, \I{PF}) von
\VonBis{60}{100}\PerMin* und wirft dabei jedes Mal ein \I{Schlagvolumen} (\I{SV}) von ca.
\VonBis{70}{80}\Ml* aus
\cite{Silbernagl3-Herz,Silbernagl3-Kreislaufsystem}. Daraus
ergibt sich \E{in Ruhe} ein \I{Herzminutenvolumen} (\I{HMV})
von ca.
\Range{5}{8}\LPerMin*.
Bei Belastung kann das Herzminutenvolumen beträchtlich
gesteigert werden. Durch die Pumpfunktion des Herzens
entsteht ein lebenswichtiger Druck im (arteriellen) Gefäßsystem, der
\I{Blutdruck}.
Dieser sorgt dafür, dass das Gewebe von Blut durchflossen
wird. Der arterielle Blutdruck lässt sich
einfach mittels Blutdruckmanschette messen. (\FullPrettyRef{topic:blutdruck})
}
{%
\begin{itemize}
  \item Herzfrequenz, Pulsfrequenz ({HF}, PF)
  \VonBis{60}{100}\PerMin*
  \item Schlagvolumen \VonBis{70}{80}\Ml*
  \item {Herzminutenvolumen} \Range{5}{8}\LPerMin*.
  \item Blutdruck:
  \begin{itemize}
    \item Druck im arteriellen Gefäßsystem
    \item Durchblutung des Gewebes
    \item Blutdruckmessung $\rightarrow$ \ABCDE{D}
  \end{itemize}
\end{itemize}
}
{}
{}
{}




\TextSlide{Einschätzung des Kreislaufs}{%
Die erste Einschätzung des Kreislaufes dauert nur wenige
Sekunden. Durch die Beurteilung der \EE{Hautfarbe} und
\EE{-temperatur}, des \EE{Radialispulses} (gut oder
schlecht tastbar, geschätzte Schnelligkeit, Rhythmik), sowie
der \EE{Rekap-Zeit} an den Fingerkuppen kann man auf
einfache Weise die Kreislauffunktion beurteilen. Bei der
Rekap-Zeit wird ein leichter Druck auf das Nagelende ausgeübt bis
das Nagelbett weiß ist, dann lässt man los. Innerhalb von
\VonBis{1}{2} Sekunden sollte das Nagelbett wieder rosarot
sein. Eine verlängerte Rekapzeit deutet auf eine
lokale Durchblutungsstörung\footnote{\E{Verlängerte
Rekapzeit}: Unter \enquote{lokale Durchblutungsstörung} fallen
auch in Folge von Kälte minderdurchblutete Extremitäten, also \enquote{kalte Finger}.}
oder Zentralisierung hin. 


}{%
\begin{itemize}
  \item Hautfarbe, -temperatur 
  \item Schweiß
  \item Rekap-Zeit 
  \item Radialispuls 
\end{itemize}
}
{}
{}
{}







\TextSlide{Blutdruck}
{\label{topic:blutdruck}%
Durch die Pumpfunktion des Herzens
entsteht ein lebenswichtiger Druck im 
(arteriellen) Gefäßsystem, der
\I{Blutdruck}.
Da das Herz rhythmisch kontrahiert, kommt es infolge des
Blutauswurfes zu synchronen -- ebenfalls rhythmischen --
Druckveränderungen. Während des Auswurfs von Blut in das
Gefäßsystem (Herzanspannung, \I{Systole}) spricht man vom
systolischen Blutdruck, während der Entspannungsphase des
Herzens (\I{Diastole}) vom diastolischen Blutdruck.
Natürlich ist der Blutdruck bei der Anspannung des Herzens
höher als bei der Entspannung (\E{Systolischer Blutdruck >
Diastolische Blutdruck)}. Zur Durchführung der
Blutdruckmessung vgl.
\FullPrettyRef{topic:rr-messung}.

\subparagraph{{Interpretation des Blutdruckwerts}}
{\tiny Werte beziehen sich auf den Erwachsenen.}
\begin{itemize}
    \item \TabSub{Normbereich}: zwischen \RRmmHg{90}{60} und
    \RRmmHg{140}{90} \index{Blutdruck!Normbereich}
    \item \TabSub{Hypotonie}: RR {\textless}~\RRmmHg{90}{60}
    \index{Hypotonie}
    \item \TabSub{Hypertonie}: RR
    {\textgreater}~\RRmmHg{140}{90}
    \index{Hypertonie}
\end{itemize}
Eine Hypertonie ist zwar eine dauerhaft behandlungsbedürftige Erkrankung, aber
nicht automatisch im Akutfall von Bedeutung. Die Patienten sind
normalerweise an ihren erhöhten Blutdruck gewöhnt und haben keine
Beschwerden. In der Regel stellen sich erst bei einer \E{plötzlichen} Erhöhung Symptome ein.
\EE{Äußere Einflüsse} haben außerdem einen großen Einfluss auf den Blutdruck, 
sodass dieser auch erhöht sein kann, ohne dass hierfür eine Erkrankung vorliegt. 
Weiters ist der untere Normalwert sehr individuell,
insbesonders bei sehr schlanken oder bettlägrigen Patienten
sind systolische Blutdruckwerte um die 90\MmHg* möglich.
Bei der Interpretation des Wertes muss daher immer bedacht
werden:
\begin{itemize}
    \item \EE{Zustand} des Patienten
    \item \EE{Schmerzen}
    \item Der \EE{sonst übliche} Blutdruck des Patienten
    (erfragen, Blutdruckprotokoll)
    \item Der \EE{Verlauf} des Blutdruckes während der
    Behandlung (Traumapatienten!)
    \item Der \EE{Erregungszustand} des Patienten (Nervosität, \ldots).
\end{itemize}

\minisec{Beispiele}

\emph{Ein 65-jähriger, männlicher Patient klagt über
Schwindelgefühl, sie messen einen Blutdruck von
\RRmmHg{100}{70}. Im Blutdruckprotokoll finden sie für die
vergangenen vier Tage folgende Einträge:
\RR{170}{110}, \RR{140}{115}, \RR{175}{113}, \RR{183}{112}, \ldots}


Wie beurteilen Sie den Blutdruck?
\index{Blutdruck!Beurteilung}

\begin{quote}
Der Blutdruck ist für den Patienten zu niedrig, der Körper
ist auf eine so plötzliche Umstellung nicht vorbereitet
gewesen. Durch das Anamnesegespräch erfahren sie, dass er
heute das erste Mal einen \enquote{Nitro-Spray} verwendet hat. Sie
vermerken dies am Einsatzprotokoll. Im Krankenhaus erfahren
Sie, dass wahrscheinlich eine Überdosierung dieses
Medikaments die Ursache dieses Blutdruckabfalls war.
${\rightarrow}$ Obwohl der Patient laut Lehrbuch einen
\enquote{schönen} Blutdruck hatte, war dieser dennoch zu niedrig!
\end{quote}

\emph{Ein 56-jähriger Patient klagt über starke kolikartige
Schmerzen im linken Unterbauch. Sie messen einen Blutdruck
von \RRmmHg{155}{95}. Der Patient kennt seinen
\enquote{gewöhnlichen} Blutdruck
nicht.}  

Wie beurteilen Sie den Blutdruck?

\begin{quote}
Der Blutdruck ist laut Lehrbuch erhöht, aber dies ist
erklärbar durch die starken Schmerzen und den damit
verbundenen Erregungszustand. Die Beschwerden stehen
wahrscheinlich in
keinem ursächlichen Zusammenhang mit dem Blutdruck.
\end{quote}

\emph{Eine 50-jährige Patientin hat Nasenbluten. Sie
messen einen Blutdruck von \RRmmHg{220}{120}. Auf Nachfrage gibt sie an
etwas Kopfsschmerzen zu haben, und etwas schwindlig sei sie auch.} 

Wie beurteilen Sie den Blutdruck? 

\begin{quote}
Der Blutdruck ist sehr stark erhöht, außerdem zeigt die
Patientin Zeichen einer hypertensiven Krise (Nasenbluten,
Kopfweh, Schwindel;
\FullPrettyRef{topic:hypertensive-krise}).
Hier ist der erhöhte Blutdruck Ursache für die Beschwerden
und muss mittels Medikamenten unbedingt gesenkt werden.
\end{quote}

}{%
\begin{itemize}
  \item 2 Phasen:
  \begin{description}
    \item[systolisch:] Maximaldruck in der Arterie am
    Höhepunkt der
    Kammerkontraktion
    \item[diastolisch:] Druck in der Arterie während die
    Kammer erschlafft ist
  \end{description}
  \item Durchführung:
  \FullPrettyRef{topic:rr-messung}
  \item Interpretation:
  \begin{itemize}
    \item {Normbereich}: zwischen \RRmmHg{90}{60} und
    \RRmmHg{140}{90}
    \item {Hypotonie}: RR {\textless}~\RRmmHg{90}{60}
    \item {Hypertonie}: RR {\textgreater}~\RRmmHg{140}{90}
    \item Bedenke:
    \begin{itemize}
      \item {Zustand} des Patienten
      \item sonst üblicher Blutdruck
      \item {Verlauf}
      \item {Erregungszustand}
    \end{itemize}
  \end{itemize}
\end{itemize}
}{}{}{}







\TextSlide
{{STU}* | Schnelle Traumauntersuchung}
{%
Liegt ein Verdacht auf ein Trauma
vor, wird eine schnelle Traumauntersuchung (STU)
durchgeführt\footnote{Wenn Verletzungen auf einen 
Körperteil beschränkt sind und
der Unfallhergang eindeutig ist (Amputation eine Fingers
durch eine Säge, \ldots) oder Bagatellverletzungen vorliegen
(eingezogener Schiefer, \ldots), können u.\,U. manche
Untersuchungen entfallen. Es muss aber bedacht werden, dass
auch kleine Verletzungen aufgrund von zugrundeliegenden
Erkrankungen entstehen, oder Folgeverletzungen vorliegen
können!}. 
Ziel ist das Aufspüren von akut
\E{lebensgefährlichen} Verletzungen (insbesonders Blutungsquellen und Schäden des ZNS). 
Sie beschränkt sich auf
die wichtigsten Organeund  besteht aus der  Inspektion
  und dem Abtasten von
  \begin{compactenum}
    \item \EE{Kopf} inkl. Ohren,
    \item \EE{Hals},
    \item \EE{Thorax} 
    (inkl. \T{Atemprobe}: Bei der Einatmung wird mit den flachen Händen
    gegen den Brustkorb gedrückt. 
    Dadurch kann ein Hinweis auf eine mechanisch-bedingte Atmungsstörung 
    gewonnen werden, z.\,B. durch eine Serienrippenfraktur),
    \item \EE{Bauch},
    \item \EE{Becken},
    \item \EE{Oberschenkel} und
    \item \EE{Rücken} (evtl. erst beim Umlagern auf eine Schaufeltrage oder Spine Board).
  \end{compactenum}
}
{\begin{itemize}
  \item Inspektion
  und Abtasten von
  \begin{compactenum}
    \item \EE{Kopf} inkl. Ohren,
    \item \EE{Hals},
    \item \EE{Thorax} (Atemprobe),
    \item \EE{Bauch},
    \item \EE{Becken},
    \item \EE{Oberschenkel},
    \item \EE{Rücken}
  \end{compactenum}
\end{itemize}}
{}
{}
{}




\TextSlide
{{\eye} Befunde}
{Erkennen kann man z.\,B. Schockzeichen,
sichtbare Blutungen, Herzrhythmusstörungen (Tachy-/Bradykardie,
Arrhythmie) sowie Veränderungen des Blutdrucks. Diese
Befunde müssen mit der Gesamtheit der Befunde, der möglichen
Diagnosen des Patienten, den Umständen und dem Zustand
des Patienten, sowie unter Berücksichtigung der Anamnese
beruteilt werden.

Bei der Schnellen Trauma-Untersuchung (STU) können
Blutungsquellen oder Hinweise auf Verletzugen des zentralen
Nervensystems (ZNS) gefunden weden. Als Nebenbefund kann bei
Verletzungen des Thorax eine mechanisch bedingte Atemstörung
festgestellt werden (\ABCDE{B}-Befund).

}{%
\begin{itemize}
  \item Sichtbare Blutungen
  \item Schockzeichen
  \item Rhythmus (Tachy-/Bradykardie, Arrhythmie)
  \item Blutdruckveränderungen
  \item STU: 
  \begin{itemize}
    \item Blutungsquellen
    \item Verletzungen des ZNS (\ABCDE{D}-Befund)
    \item Mechanische Atemstörungen (\ABCDE{B}-Befund)
  \end{itemize}
\end{itemize}
}
{}
{}
{}

% 
% \begin{PictureSeriesWide}{}{Bilderserie ABCDE}
% \PictureSeriesRowWide{3}
% {./images/trauma/pupillendifferenz.jpg}
% {Pupillendifferenz}
% {./images/motal-aass/00800/neuro6.jpg}
% {BZ-Messung \SourcePicture{Motal}}
% {./images/hirtler-aass/Riva-Rocci_edited.pdf}
% {Prinzip
%     der auskultatorischen RR-Messung: In Abhängigkeit vom
%     außen anliegenden Manschettendruck kommt es zu einem
%     unterschiedlich starken Blutfluss und dadurch
%     entstehenden, hörbaren, Strömungsgeräuschen.}
% {empty}
% {}
% \end{PictureSeriesWide}

\Text{Typische Maßnahmen}
{%
\begin{itemize}
  \item Lagerung mit erhöhten Beinen
  \item Blutstillung
  \item \TxMassVit
\end{itemize}
}

%%%%%%%%%%%%%%%%%%%%%%%%%%%%%%%%%%%%%%%%%%%%%%%%%%%%%%%%%%%
%%%%%%%%%%%%%%%%%%%%%%%%%%%%%%%%%%%%%%%%%%%%%%%%%%%%%%%%%%%
%%%%%%%%%%%%%%%%%%%%%%%%%%%%%%%%%%%%%%%%%%%%%%%%%%%%%%%%%%%
%%%%%%%%%%%%%%%%%%%%%%%%%%%%%%%%%%%%%%%%%%%%%%%%%%%%%%%%%%%
\clearpage
\subsection{{\FontKeyboardFamily{D}} Neurologisch
(Disability)}

\TextSlide{Cute Eyes}{%
Es wird die \EE{Be\-wusst\-seins\-quali\-tät} (Orientierung zur Person,
Raum, Zeit und Situation) oder, wenn angemessen, alternativ 
die Glasgow Coma Scale (\EE{GCS}) bestimmt. Anschließend werden die
\EE{Pupillen} und deren Lichtreaktion untersucht. Die
\EE{grobe Motorik} wird an der oberen Extremität mittels
Händedruck (Kraftprobe) überprüft, wenn angemessen wird auch 
die untere Extremität geprüft. Wenn der Verdacht auf neurologische 
Ausfälle besteht, muss überprüft werden ob der
    Patient \EE{Hände} und \EE{Füße} spüren und aktiv
    bewegen kann.
Wenn erforderlich (neurologische Defizite, Hinweise auf Stoffwechselstörungen, \ldots) 
erfolgt eine \EE{Blutzuckermessung}.

}{
\begin{itemize}
  \item Vitalwerte (\EE{HF}, \EE{RR})
  \item Neurologisch:
  \begin{itemize}
    \item \EE{Orientierung} oder *\EE{GCS}
    \item \EE{Pupillen}
    \item \EE{Kraftprobe OE}, *UE
    \item Ggfs.: Kann der
    Patient \EE{Hände} und \EE{Füße} spüren und aktiv
    bewegen?
  \end{itemize}
  \item Ggfs. \EE{Blutzuckermessung}
\end{itemize}
}{}{}{}



\TextSlide{Orientierung}
{%
\label{topic:orientierung}
\index{Orientierung}
%
Es ist wichtig zu Überprüfen ob der Patient orientiert ist.
Man unterscheidet dabei die Orientierung zu:  
\begin{itemize}
  \item \EE{Person}: Name, Alter.
  \item \EE{Zeit}: Tageszeit,
  Wochentag, Monat, Jahreszeit, Jahr?
  \item \EE{Ort}: Wo befindet sich der Patient gerade?
  \newline
  (Stadt, Land, Adresse, Wohnung, Auto, etc.)
  \item \EE{Situation}: Was passiert gerade?
\end{itemize}

}{
\begin{itemize}
    \item Zur Person
    \item Zeitlich
    \item Örtlich
    \item Zur Situation
\end{itemize}
}{}{}{}



%%%%%%%%%%%%%%%%%%%%%%%%%%%%%%%%%%%%%%%%%%%%%%%%%%%%%%%%%%%
\TextSlide{Glasgow Coma Scale (GCS)}
{%
\index{Glasgow Coma Scale}
\index{GCS}
\label{topic:gcs}
%
Die Glasgow Coma Scale ist eine spezielle Skala, welche der
objektiven Einschätzung der Bewusstseinslage  dient. Sie
hat einen Wertebereich von \VonBis{3}{15} Punkten. Ein voll
bewusstseinsklarer Patient hat 15 Punkte, ein schwerst
bewusstseinsgestörter oder toter Patient 3 Punkte (nicht
0!). Bewertet werden die Augenöffnung, die beste verbale und
die beste motorische Reaktion des
Patienten.\ZFootnote{Regelmäßig findet man
Handlungsanweisungen -- überwiegend für ärztliches Personal
-- in der Literatur, welche die \Abk{GCS} als
Entscheidungskriterium (z.\,B. für
die Intubation) heranziehen. Der Nutzen ist aber genauso
regelmäßig umstritten.}
}{
\VonBis{3}{15} Punkte.
\begin{itemize}
    \item Augenöffnung (4)
    \item Beste verbale Reaktion (5)
    \item Beste motorische Reaktion (6)
\end{itemize}
\prettyref{tab:gcs}.
}{}{}{}

\TextSlide
{{\eye} Befunde}
{Erkennen lassen sich z.\,B.  Desorientierung, pathologische
Pupillenreaktionen (abnorme Weitung oder Engstellung,
Differenz, verlangsamte oder fehlende Reaktion (Fixierung),
Entrundung), motorische (Lähmungen) oder sensible (Kribbeln,
Gefühlsverlust) Ausfälle sowie Zuckerstoffwechselstörungen.
}
{
\begin{itemize}
  \item Entgleiste Vitalwerte
  \item Desorientierung
  \item Pupillen: Abnorme \E{Weite}, \E{Differenz},
  \E{Reaktion}, \E{Entrundung}
  \item Motorische und sensible Ausfälle
  \item Zuckerstoffwechselstörungen
\end{itemize}
}
{}
{}
{}

\begin{PictureSeriesWide}{}{Bilderserie ABCDE}
\PictureSeriesRowWide{3}
{./images/trauma/pupillendifferenz.jpg}
{Pupillendifferenz \SourcePicture{GABS}}
{./images/motal-aass/00800/neuro2.jpg}
{Kraftprobe an den Armen \SourcePicture{Motal}}
{./images/motal-aass/00800/neuro6.jpg}
{BZ-Messung \SourcePicture{Motal}}
{empty}
{}
\end{PictureSeriesWide}
% 
% \Text{Typische Maßnahmen}
% {%
% \begin{itemize}
%   \item Strategie- und Zielentscheidung
% \end{itemize}
% }

%%%%%%%%%%%%%%%%%%%%%%%%%%%%%%%%%%%%%%%%%%%%%%%%%%%%%%%%%%%
%%%%%%%%%%%%%%%%%%%%%%%%%%%%%%%%%%%%%%%%%%%%%%%%%%%%%%%%%%%
%%%%%%%%%%%%%%%%%%%%%%%%%%%%%%%%%%%%%%%%%%%%%%%%%%%%%%%%%%%
%%%%%%%%%%%%%%%%%%%%%%%%%%%%%%%%%%%%%%%%%%%%%%%%%%%%%%%%%%%
\clearpage
\subsection{{\FontKeyboardFamily{E}} Erweiterte Untersuchung}

\TextSlide{E -- Wenn es nicht so klar ist: Raum für Instinkt und Erfahrung}{%
\ZFootnote{E: \Lang{engl.} \E{Exposure}.} 
Mit den Punkten \Range{\ABCDE{A}}{\ABCDE{D}}
 wären die wichtigsten Untersuchungen
durchgeführt.
Dennoch gibt es einige Patientengruppen, die zwar bei den
ABCDs nicht auffällig sind, aber dennoch als vital bedroht
(kritisch) angesehen werden müssen. 
Diesem Umstand trägt der Punkt E Rechnung: 
Abhängig von den Leitsymptomen, dem
bisher festgestellten Zustand des Patienten und dem Wissen des
Fachpersonals sollen zielgerichtete Untersuchungen durchgeführt oder
Anamanesefragen erhoben werden, um eine vitale Bedrohung zu bestätigen
oder als unwahrscheinlich einzustufen. 
Der Punkt E kann und soll \E{nicht} standardisiert sein,
sein Erfolg ist definitiv von der Erfahrung und von dem Wissen des Untersuchers
abhängig. Im Punkt E können sämtliche Untersuchungen und
Anamneseschritte durchgeführt werden, welche zur Festellung (oder
Quasi-Ausschluss) der vitalen Bedrohung notwendig sind.

\begin{quote}
  Ein Beispiel für so eine Patientengruppe wäre der Patient mit
        wiederholten Krampfanfällen, der sich als St.p. Kollaps-Patient
        präsentiert (Schlüsselfragen: \Speech{Hat der
        Patient gekrampft? Wie oft hat er gekrampf?}).
\end{quote}

}{%
\begin{itemize}
  \item Wenn Einschätzung der vitalen
  Bedrohung oder Hauptbeschwerde unklar ist.
  \item Wichtig:
  \begin{itemize}
    \item Erfahrung
    \item Hintergrundwissen
  \end{itemize}
\end{itemize}
}{}{}{}

\TextSlide{Das heißt?}{%
Ist die vitale
Bedrohung also noch nicht klar einzuschätzen, müssen
\EE{gezielt weitere Untersuchungen oder Anamneseschritte} durchgeführt
werden. Dies beinhaltet auch in der Situation besonders
relevante Teile der Anamnese, u.\,U. auch eine
\E{Fremdanmnese}. 
Dafür ist es erforderlich, sich \E{zuerst
die Situation, die bekannten Beschwerden und die bereits
erhobenen Befunde zu vergegenwärtigen und
daraus entsprechende Schlüsse und \EE{Verdachtsmomente}}
abzuleiten. 

Es müssen dann zielgerichtete Fragen oder
Untersuchungen zur Abklärung von \E{Alarmzeichen}, insbesonders Alarmdiagnosen und relevante Begleiterkrankungen gestellt bzw. durchgeführt werden.
Begleiterkankungen können sich erschwerend auf die
Beschwerden des Patienten auswirken (z.\,B.
vorgeschädigetes, schwaches Herz bei vorbestehender
Herzinsuffizienz, \ldots). Zwei Erkrankungen, welche für
sich nicht unbedingt besonders gefährlich sind,
können in \EE{Kombination} gefährlich sein. 

Besonders bei Trauma-Patienten kann es notwendig sein, den
Patienten zu \EE{entkleiden} um beispielsweise
\E{Prellmarken} zu erkennen. Das Entkleiden des Patienten
muss unbedingt unter \Ca{Wahrung des Wärmeerhalts} erfolgen
und darf i.\,d.\,R. nur in warmer Umgebung und unter Schutz
vor Witterungseinflüssen erfolgen. Eventuell muss dieser
Schritt so lange verzögert werden bis sich der Patient in
einer warmen Umgebung befindet.

Bei Punkt E ist ein solides \EE{Hintergrundwissen} über
relevante Erkrankungen, Verletzungen entscheidend, je
souveräner das Wissen beherrscht wird, desto gezielter
können Fragen und Befunde erhoben werden und desto besser
wird die Einschätzung des Patienten erfolgen.

}{%
\begin{itemize}
  \item Zuerst: \EE{Situation und bereits
erhobenen Befunde vergegenwärtigen und
daraus entsprechende Schlüsse und {Verdachtsmomente}
ableiten}
  \item Zielgerichtete Fragen auf Alarmzeichen, insb.
  Alarm-Diagnosen und rel. Begleiterkrankungen
  \item Inkl. relevante Teile der (Fremd-)Anamnese
  \item Ggfs. entkleiden (in warmer Umgebung)
\end{itemize}
}{}{}{}

\TextSlide
{{\eye} Befunde}
{Zu erkennen sind sonstige Hinweise auf
eine vitale Bedrohung, welche bisher nicht erkennbar waren.
So kann Beispielsweise ein \E{Augenzeuge} auf Nachfrage
berichten, dass der aufgefundene Patient in seinem Beisein
zweimal gekrampft hat, was ein Hinweis für einen
beginnenden Status epilepticus (und damit eine vitale
Bedrohung) sein kann. \E{Prellmarken} können Hinweise auf
innere Verletzungen sein (zweizeitige Milzruptur,
\ref{topic:milzruptur-zweizeitig}). \E{Einstiche} können ein
Hinweis auf eine Intoxikation sein (wobei diese nicht die
Ursache des aktuellen Problem sein muss!).
}
{
\begin{itemize}
  \item Sonst. Hinweise auf vitale Bedrohung
  \item Anamnese (z.\,B.
  $\geq$ 2 Krämpfe/24\Hour*)
  \item Gefährliche Kombinationen von
  Erkankungen
  \item Prellmarken
  \item \ldots 
\end{itemize}
}
{}
{}
{}
%%%%%%%%%%%%%%%%%%%%%%%%%%%%%%%%%%%%%%%%%%%%%%%%%%%%%%%%%%%
%%%%%%%%%%%%%%%%%%%%%%%%%%%%%%%%%%%%%%%%%%%%%%%%%%%%%%%%%%%
%%%%%%%%%%%%%%%%%%%%%%%%%%%%%%%%%%%%%%%%%%%%%%%%%%%%%%%%%%%
%%%%%%%%%%%%%%%%%%%%%%%%%%%%%%%%%%%%%%%%%%%%%%%%%%%%%%%%%%%
\clearpage
\subsection{{\FontKeyboardFamily{=}} Beurteilung}
\TextSlide{Karten auf den Tisch!}{%
Die \EE{Einschätzung der vitalen Bedrohung}
erfolgt \E{ständig} während des gesamten
Einschätzungsblocks. Ein Patient kann jederzeit anhand
der Alarmzeichen für \E{vital bedroht} erklärt werden. 

Im Verlauf des Einschätzungsblocks gewinnt man erste
Hinweise auf mögliche Diagnosen und kann meist eine
vorläufige \E{Verdachtsdiagnose} formulieren (\Speech{\enquote{Was
ist das Problem des Patienten?}}).

\EE{Sofortmaßnahmen} müssen bei Bedarf sofort ergriffen
werden und unterbrechen evtl. den Einschätzungsblock. Sie
können jedoch auch von einer zweiten Person parallel
durchgeführt werden. Wenn die jeweiligen Sofortmaßnahmen durchgeführt
wurden, wird der Block fortgesetzt. Typische Sofortmaßnahmen
sind die Maßnahmen des Selbst- oder Fremdschutzes, eine
situationsgerechte Lagerung oder die \EE{\TxMassVit}.
Vgl. \prettyref{topic:patientenmanagement-sofortmassnahmen}.
}{%
\begin{itemize}
  \item \EE{Einschätzen der vitalen Bedrohung}
  ständig während des gesamten Blocks
  \item Vorläufige Verdachtsdiagnose formulieren
  \item[\eye] Alarmzeichen
  \item[\eye] Hinweise auf mögliche Diagnosen
\end{itemize}
}{}{}{}



\subsection{Zeichen einer Vitalen Bedrohung: Alarmzeichen (\protect\enquote*{Red
Flags} \RedFlag)}

\TextSlide{Alarmzeichen (\protect\enquote*{Red
Flags} \RedFlag)}
{%
Das rasche Feststellen einer vitalen Bedrohung ist besonders
wichtig. Die vitale Bedrohung kann sich aus drei wesentlichen
Säulen ergeben: \marginpar{\includegraphics[width=\linewidth]{./images/teaser/minen_hinweis2.jpg}}

\begin{itemize}
  \item[\RedFlag] \EE{Massive \T[3ABC-Störung]{\ABCDE{3}\ABCDE{A}\ABCDE{B}\ABCDE{C}-Störung}}: Massive
  Störungen von Vitalfunktionen 1. Ordnung (Bewußtsein \ABCDE{3}, Atmung inkl. Atemweg \ABCDE{A}\ABCDE{B},
  Kreislauf \ABCDE{C})
  \item[\RedFlag] \T{Alarmsymptome}: Zu den Alarmsymptomen
  gehören u.\,a.
  Atemnot, die sich nicht bessert, brodelndes Atemgeräusch,
  Thoraxschmerz,
  Schocksymptome,
  entgleiste Vitalwerte,
  schwere Verletzungen,
  Hirndruckzeichen, etc.
  \item[\RedFlag] \T{Alarmdiagnosen}: Manche Diagnosen
  sind automatisch mit einer vitalen Bedrohung gleichzusetzen,
  auch wenn eventuell Symptome fehlen oder nur schwach
  ausgeprägt sind. Dazu gehören beispielsweise der
  Herzinfarkt, das kardiale oder toxische Lungenödem,
  der Status Epilepticus, eine Beckenfraktur, u.\,s.\,w.
\end{itemize}

% \parpic[r][4cm]{\includegraphics[width=4cm]{./images/teaser/minen_hinweis2.jpg}}


 Während des weiteren Behandlungsverlaufs
kann man auf weitere Befunde oder Fakten stoßen,
welche auf eine vitale Bedrohung hinweisen.

}{\vspace{1em}

\fcolorbox{DefaultRed}{White}{
\begin{minipage}[t]{\linewidth - 4\Korrektur}
{
\begin{itemize}
    \item[\RedFlag] Massive \ABCDE{3}\ABCDE{A}\ABCDE{B}\ABCDE{C}-Störung
    \item[\RedFlag] Alarm-Symptome
    \item[\RedFlag] Alarm-Diagnosen
\end{itemize}

Übersicht \FullPrettyRef{tab:alarmzeichen}




}
\end{minipage}
}
}
% {./images/teaser/minen_hinweis2.jpg}
% {Vorsicht!}
% {GABS.}

\Cave{Sobald eine vitale Bedrohung erkannt wird, müssen die
\TxMassVit{}
\E{frühzeitig}
ergriffen werden (evtl. sogar noch vor Abschluss des
Einschätzungsblockes)!}



\subsection{Sofortmaßnahmen}
\label{topic:patientenmanagement-erstmassnahmen}
\label{topic:patientenmanagement-sofortmassnahmen}


Sofortmaßnahmen müssen bei Bedarf sofort ergriffen
werden und unterbrechen evtl. den Einschätzungsblock. Sie
können jedoch auch von einer zweiten Person parallel
durchgeführt werden. Wenn die jeweiligen Sofortmaßnahmen durchgeführt
wurden, wird der Block fortgesetzt. Typische Sofortmaßnahmen
sind z.\,B. 
\begin{itemize}
  \item Maßnahmen des Selbst- oder Fremdschutzes (Retten
  aus dem Gefahrenbereich, GAS-Maßnahmen, \ldots),
  \item Einschätzung der Indikation zur Wirbelsäulenimmobilisation
  \item eine situationsgerechte Lagerung,
  \item die \EE{\TxMassVit} oder auch die
  \item Vorgezogene Transportentscheidung bei absolut zeitkritischen Patienten.
\end{itemize}

% M %%%%%%%%%%%%%%%%%%%%%%%%%%%%%%%%%%%%%%%%%%%%%%%%%%% M %
% Einschätzung der Indikation zur Wirbelsäulenimmobilisation
\ProcedurePrintDefault{YY14200B}
\Commentary[Konsensus]{Review 2011-01-11}

% \TextSlide{Vital bedrohte Patienten}{%
% Bei vital bedrohten Patienten gibt es 
% \EE{fünf \enquote{{\TxMassVit}}},
% die \E{sobald wie möglich} (evtl. sogar
% noch vor Abschluss des Ersteinschätzungsblockes) zu erfolgen
% haben (\begin{inparaenum}
% \item Situationsgerechte Lagerung, 
% \item Sauerstoffgabe (sofern keine Kontraindikation), 
% \item Notarztnachforderung (mit Ausnahmen), 
% \item Reanimationsbereitschaft, 
% \item Engmaschiges, bestmögliches Monitoring
% \end{inparaenum}
% ).
% }{
% \begin{itemize}
%   \item {\TxMassVit}
% \end{itemize}
% }{}{}{}



% M %%%%%%%%%%%%%%%%%%%%%%%%%%%%%%%%%%%%%%%%%%%%%%%%%%% M %
% Standardmaßnahmen bei vital bedrohten Patienten
\ProcedurePrintDefault{YY13100B}


% M %%%%%%%%%%%%%%%%%%%%%%%%%%%%%%%%%%%%%%%%%%%%%%%%%%% M %
% Reanimationsbereitschaft
\ProcedurePrintDefault{YY13210B}


\TextSlide
{\underline{Sonderfall:} Vorgezogene Transportentscheidung bei absolut zeitkritischen Patienten}
{%
\label{topic:sonderfall-absolut-zeitkritische-patienten}
\index{Vorgezogene Transportentscheidung
bei absolut zeitkritischen Patienten}
Bei \Ca{absolut zeitkritischen Patienten}, bei denen bereits
die ABC-Einschätzung ergibt, dass ein Transport nicht
aufschiebbar ist, kann schon \Ca{nach dem Punkt C eine
Transportentscheidung} (Spital,
Notarzt-Rendez-vous, \ldots, vgl.
\FullPrettyRef{topic:pam-strategie}) getroffen werden.
Die Maßnahmen von D und E sollen dann, sofern möglich,
während des Transportes erfolgen.
Diese Einschätzung setzt jedoch ein hohes Maß an Erfahrung
voraus und ist auch
\EE{von regionalen Gegebenheiten abhängig!}


Beispiele für ein solches Vorgehen wären z.\,B. nicht
beherrschbare starke Blutungen oder geburtsunmögliche Lagen.
}
{
\begin{itemize}
  \item \Ca{Im Ausnahmefall schon nach dem Punkt C Strategie- und Zielentscheidung}
  \item Z.\,B. nicht beherrschbare starke 
  Blutungen, geburtsunmögliche Lagen
\end{itemize}
}
{}
{}
{}

\Text{Nicht vital bedrohte Patienten}{%
Auch bei
nicht vital bedrohten Patienten können Sofortmaßnahmen
ergriffen werden (z.\,B. Lagerung). Es muss individuell
überlegt werden, welche Maßnahmen zuerst zu treffen sind
und welche warten können. 
}{

}{}{}{}




\subsection{Strategieentscheidung}

\label{topic:pam-strategie}
\label{topic:strategieentscheidung}

Während und nach dem Einschätzungsblock (und den Sofortmaßnahmen) müssen
die Prioritäten des Patienten gesetzt und der weitere
Plan erstellt werden. Daraus ergeben sich die weiteren
Handlungsschritte. Ziel ist es, zu erkennen, welche
Maßnahmen, Untersuchungen und Fragen am dringlichsten sind
und zuerst durchgeführt werden müssen, und was eher in der
Abfolge nach hinten ver\-schoben wird; d.\,h. der weitere
Ablauf ist situations- und patientenabhängig.

\begin{itemize}
  \item \Speech{Besteht eine hohe Transportpriorität?}
  \item \Speech{Sind noch weitere Untersuchungen
  oder Fragen zur Lagebeurteilung nötig? (z.\,B.
  Anzahl der Krampfanfälle, \ldots )}
  \item \Speech{Welche weiteren Untersuchungen, Maßnahmen
  und Fragen sind wichtig?}
  \item \Speech{Was ist wichtig, was kann warten?}
\end{itemize}


\TextSlide
{Transportpriorität}
{%
\index{Transportpriorität}%
\label{topic:transportprioritaet}%
Bei manchen Krankheitsbildern bringt eine ausgiebige
Behandlung vor Ort keinen Vorteil, andere hingegen
profitieren stark von erweiterten Maßnahmen \E{vor} dem
Transport. Dementsprechend muss überlegt werden, ob weiter
vor Ort untersucht oder behandelt werden soll, oder ob ein
rascher Transport in der jeweiligen Situation sinvoller ist.
Zu bedenken ist hierbei:

\begin{enumerate}
  \item \EE{Zeitkritisch?} Ergibt sich aus der
  Verdachtsdiagnose, z.\,B. ein Schlaganfall ist sehr
  zeitkritisch
  \item \EE{Stabilisierung notwendig?} Ergibt sich ebenfalls
  aus der Verdachtsdiagnose bzw. dem Patientenzustand, z.\,B.
  ein brodelndes Lungenödem muss zuerst vor Ort behandelt und stabilisiert werden
  bevor der Patient transportfähig ist.
  \item Ggfs. \EE{Verfügbarkeit eines Notarztmittels}: Wenn
  ein Notarztmittel notwendig sein sollte (z.\,B. Patient hat
  Alarmzeichen, \ldots) muss die erforderliche Zeitspanne
  bis zum Eintreffen bedacht werden (diese ist ggfs. über
  die Leitstelle zu erfragen). Die Verfügbarkeit ist
  abzuwägen gegen die \ldots
  \item \EE{Transportdauer zur Zieleinrichtung}: Ist die
  erwartete Transportdauer kürzer als die Eintreffzeit eines
  Notarztmittels ist -- sofern der Patient nicht stabilisert
  werden muss -- der Transport oft vorteilhafter. 
  \item \EE{Rendez-vous}: Beim Rendez-vous (\Lang{franz.}
  Treffen) treffen sich das Rettungsmittel und das
  Notarztmittel auf der Strecke, sodass einerseits ein
  schneller Transport, und andererseits eine ärztliche
  Versorung ohne große Zeitverzögerung möglich ist. Es
  besteht jedoch ein gewisser Organisationsaufwand, welcher
  sich erst bei größeren Distanzen auszahlt. Daher wird
  diese Strategie eher in ländlichen, und kaum in
  städtischen Gebieten angewendet.
\end{enumerate}

Wie unter \FullPrettyRef{topic:sonderfall-absolut-zeitkritische-patienten},
bereits ausgeführt, kann als Sofortmaßnahme bei 
\EE{absolut zeitkritischen Patienten}, bei denen bereits die 
ABC-Einschätzung ergibt dass ein Transport nicht aufschiebbar ist, 
schon \EE{nach dem Punkt C eine Zielentscheidung} (Spital, 
Notarzt-Rendez-vous, \ldots) getroffen werden. 
Die Maßnahmen von D und E sollen dann, sofern möglich, während des 
Transportes erfolgen. 
Darunter fallen z.\,B. nicht beherrschbare starke 
Blutungen oder geburtsunmögliche Lagen.

}
{
\begin{enumerate}
  \item \EE{Zeitkritisch?} 
  \item \EE{Stabilisierung notwendig?} 
  \item Ggfs. \EE{Verfügbarkeit eines Notarztmittels}
  \item \EE{Transportdauer} zur Zieleinrichtung
  \item Evtl. \EE{Rendez-vous}?
\end{enumerate}
}
{}
{}
{}


\subsection{Re-Evaluation (Verlaufskontrolle)}
\index{Verlaufskontrolle}

\Dictum[The bowl of petunias.\\
In: Douglas Adams: The
Hitchhiker's Guide to the Galaxy]
{Oh no, not again.}

\noindent
Der Einschätzungsblock inkl. der Strategieentscheidung muss
\EE{regelmäßig während der Betreuung} bzw.
während des Transports -- in
angemessenem Umfang -- \EE{wiederholt werden}, um eine
Verschlechterung des Patientenzustandes rechtzeitig zu erkennen. 
Änderungen in der Behandlung und der Strategie können sich
jederzeit ergeben. \Commentary[Einschätzungsblock:
Regelmäßge, angepasste Wiederholung]{%
Die Forderung nach regelmäßiger, angepasster Wiederholung
des Einschätzungsblockes entspricht dem früher
gebräuchlichen Begriff \enquote{ständige Verlaufskontrolle}. Die
Formulierungen \enquote{regelmäßig} und \enquote{in angemessenem Umfang}
sind bewusst offen gewählt worden, da dies sowohl z.\,B. für
den kritischen Bewusstlosen, als auch für den
\enquote*{munter vor sich hinplaudernden} KTW-Patienten
gelten soll.
} % Commentary

\section{Transportentscheidung}

Nach dem Einschätzungsblock erfolgt die Entscheidung über
den Transport. Diese ist abhängig
\begin{enumerate}
  \item vom Ergebnis
  des \EE{Einschätzungsblockes},
  \item der \EE{Transportpriorität}
  (\FullPrettyRef{topic:transportprioritaet}),
  \item und der resultierenden
  \EE{Strategieentscheidung}
  (\FullPrettyRef{topic:strategieentscheidung}).
\end{enumerate}


Ist ein rascher Transport notwendig, so werden rasch 
die entsprechenden Vorbereitungen getroffen
und dieser der Situation entsprechend durchgeführt. 
Ist ein rascher Transport nicht notwendig, so
können eine ausführliche Anamneseerhebung und ausführliche
Untersuchen durchgeführt werden. Dadurch kann oft eine bessere 
Verdachtsdiagnose erarbeitet werden und die Versorgung sowie das 
Transportziel optimiert werden.

\subsection{Zielentscheidung}

\parpic[r]{\includegraphics[width=0.5\linewidth]{./images/gabriel-sebastian-aass/Einfahrt_Leistsystem_IMG_1907_2_00800.jpg}}
\noindent Je nach Art des Patientenkontakts
(Rettungseinsatz, Krankentransport, Ambulanzdienst, etc.) hat die
Zielentscheidung eine unterschiedliche Bedeutung.


Weiters hat der \EE{Zustand des Patienten} sowie die
\EE{Verdachtsdiagnose} Einfluss auf den \EE{Zeitpunkt} des
Transportes (schnelle Hospitalisierung vs. Stabilisierung vor Ort).

\Text{Zielentscheidung beim Primäreinsatz}{%
Beim Primäreinsatz (\enquote{Rettungseinsätze}) ist der
Transport des Patienten die Regel. Besonders wichtig ist die
\EE{Auswahl des richtigen Zieles}. Hierbei muss bedacht werden:
\begin{itemize}
  \item Benötigte Fachabteilung
  \item Eventuell benötigte Fachabteilungen, 
  welche vor Ort (\enquote*{im Hintergrund}) verfügbar sein sollten 
  (z.\,B. Gynäkologie bei Patientinnen mit Bauchbeschwerden)
  \item Regionale Regelungen (Bettenkontingente, 
  spezielle Zuständigkeiten für bestimmte Krankheitsbilder oder 
  Patientengruppen, \ldots)
\end{itemize}
Bei kritischen Patienten
hat diese Auswahl oft große Folgen für den Erfolg der
weiteren Behandlung. Wird der Patient einer ungeeigneten
Einrichtung (unzuständige Fachabteilung, benötigte
Fachabteilung oder Spezialeinrichtungen nicht vor Ort verfügbar, \ldots) zugeführt,
kann es zu \EE{Verzögerungen} in der Behandlung kommen,
welche (vermeidbare) \EE{bleibende Schäden} verursachen
können.
}{}{}{}{}


\Text{Zielentscheidung beim Sekundäreinsatz}{%
Beim Sekundäreinsatz gibt es i.\,d.\,R. einen
vordefinierten Zielpunkt. Hier muss vor allem bedacht
werden, ob das vorhandene \EE{Material ausreicht}, um den
Patienten für die Dauer des Transportes wie gewünscht zu
versorgen. Dazu ist die Kenntnis über die speziellen
Bedürfnisse des Patienten notwendig (Medikamente,
Sauerstoff, Dekubitusprophylaxe,
\ldots).
}{}{}{}{}


\Text{Zielentscheidung bei Ambulanzdiensten}{%
\index{Triage!Patientenmanagement}%
Bei Ambulanzdiensten wird der Sanitätsdienst häufig
von Patienten in Anspruch genommen, welche keine
schwerwiegenden Gesundheitsstörungen aufweisen. Außerdem hat
bei Ambulanzdiensten die \E{Behandlung vor Ort} einen
größeren Stellenwert. Je nach materieller und
personeller Ausstattung ist oft eine definitve Behandlung
möglich.

% \parpic[r]{\includegraphics[width=4cm]{./images/TAE/erste-hilfe-dif-schild-stange-2.jpg}}
Dementsprechend werden Ambulanzpatienten oft nicht in eine
andere Gesundheitseinrichtung transportiert, sondern begeben
sich auf eigenen Wunsch selbstständig in eine weiterführende
Betreuung, oder werden nach ärztlicher Begutachtung aus der
Behandlung entlassen. Reichen die diagnostischen oder
therapeuthischen Hilfsmittel vor Ort jedoch nicht aus, oder
ist der Patient nicht in der Lage, sich selbstständig in
ärztliche Behandlung zu begeben, bzw. liegt eine
schwerwiegende Erkrankung oder Verletzung vor, so gilt das
oben beim Primäreinsatz Gesagte: Auch hier ist die Wahl der
richtigen \E{Zieleinrichtung} entscheidend. Die zusätzliche
Schwierigkeit bei Ambulanzdiensten besteht darin, jene
Patienten (rechtzeitig) zu erkennen, welche einen
Abtransport benötigen. Zusätzlich kann es notwendig sein,
Prioritäten zu setzen, wenn die verfügbare Anzahl der
Transportmittel nicht ausreicht. Bei vielen
Großambulanzdiensten werden Patienten entsprechend dem
Patientenleitsystem in die Behandlungsprioritäten
\VonBis{I}{IV} sowie die Transportprioritäten A und B
eingeteilt (\E{Triage}, \prettyref{topic:triage}).
}{}{}{}{}

\TextSlide{Bettenkontingent}{%
\index{Bettenkontingent}%
In vielen Rett\-ungs\-dienst\-be\-reichen gibt es sog.
Bettenkontingente. Durch dieses System versucht man, die
Arbeitslast möglichst gleichmäßig und gerecht auf die
einzelnen Spitäler aufzuteilen und damit Behandlungs- und
Wartezeiten zu minimieren. Die Zuteilung der Betten
(Behandlungsplätze) erfolgt zentral über eine Leitstelle
bzw. über Datenterminals. 
}{
\begin{itemize}
  \item Möglichst gleichmäßige Auslastung der
  Be\-hand\-lungs\-ein\-richt\-ung und Warte\-zeit\-re\-duzier\-ung
  \item Geeignete Abteilung
  \item Zuteilung durch eine Leitstelle
\end{itemize}
}{}{}{}

\TextSlide{\enquote{Spezialbetten}}{%
\index{Spezialbetten}%
Über ein Bettenkontingent sind i.\,d.\,R. Behandlungsplätze
von häufig benötigten Spezialabteilungen verfügbar.
Zusätzlich können jedoch spezielle Einrichtungen
erforderlich werden, dazu zählen:

\begin{itemize}
  \item \II[Ueberwachungsbett@Überwachungsbett]{Überwachungsbett} / \II{Intensivbett} für
  kritisch kranke Patienten, die einer intensivmedizinischen
  Behandlung bedürfen
  \item \II{Stroke Unit}: Behandlungseinheit für
  Schlaganfallpatienten
  \item \II{Herzkatheterlabor} (PTCA\footnote{PTCA:
  Perkutane Transluminale Coronar-Angioplastie},
  PCI\footnote{PCI: Percutane Coronar-Intervention}) für
  Herzinfarktpatienten
  \item \II{Schockraum}: Behandlungseinheit mit Vorrichtungen 
  zur Stabilisierung und Aufrechterhaltung von Vitalfunktionen,
  normalerweise von schwer
  verletzten, je nach regionaler Regelung eventuell auch von schwer 
  erkrankten Patienten. Hier wird der Patient vor einer weiteren
  Behandlung stabilisiert. Der Schockraum ist normalerweise 
  einer Abteilung für Unfallchirurgie zugeordnet.
  \item Gebärende 
  \item Psychiatrische Patienten
\end{itemize}

Bei \EE{gebärenden Patientinnen} ist meistens die
Anmeldung zur Geburt bei einer bestimmten Einrichtung entscheidend. Für
\EE{psychiatrische Patienten} gibt es oft regionale,
rechtlich begründete, Vorschriften. \opt{REGVIE}{In Wien werden
sie gemäß ihrer
\E{Meldeadresse} einer Einrichtung zugeordnet.} 
}{
\begin{itemize}
  \item {Überwachungsbett} / {Intensivbett} 
  \item {Stroke Unit}
  \item {Herzkatheterlabor} (PTCA,
  PCI) 
  \item {Schockraum}
  \item Speziell: Gebärende, psychiatrische Patienten
\end{itemize}
}{}{}{}

\Text{Sonderfall: Kurz zurückliegende Spitalsentlassung}{%
\index{Bumerang-Patient|Z}%
In manchen Rettungsdienstbereichen wurden Regelungen
eingeführt, wonach Patienten, welche innerhalb einer
gewissen Zeitspanne (z.\,B. 72\Hour*) aus einer
Spitalsbehandlung entlassen wurden, bei einer neuerlichen
Spitalsbedürftigkeit wieder in das gleiche Spital
aufgenommen werden sollen oder müssen. Das soll einerseits
einen reibungsfreieren Ablauf ermöglichen, da der Patient im
Spital bereits bekannt ist, und andererseits ein
\enquote{Abschieben} von Patienten seitens der Krankenanstalten
verhindern \Z{(\enquote{Bumerang-Patient})}.

Diese Regelung trifft natürlich nur zu, wenn das betreffende
Spital über die für den Patienten notwendigen Einrichtungen
und Abteilungen verfügt. Ist der neuerliche Grund einer
Spitalsbehandlung ein grundsätzlich anderer als der
vorherige, und kann ist das in Frage kommende Spital für das
neue Problem ungeeignet, muss der Patient natürlich in eine
\E{geeignete Einrichtung transportiert werden.}
}{}{}{}{}


\Text{Sonderfall: Bettenzusage}{%
\index{Bettenzusage}%
Spitäler können grundsätzlich auf Anfrage von sich aus sog.
\enquote{Bettenzusagen} erteilen. In der Regel erfolgt dies häufig
bei Privatspitälern und Patienten mit einer privaten
Zusatzversicherung, sowie bei speziellen Einrichtungen,
welche chronisch kranke Patienten über längere Zeit
betreuen, und in denen die Patienten bekannt sind.
Derartige Zusagen sind unabhängig von einem eventuell vorhandenem
Bettenkontingentsystem. 

Das Fachpersonal muss entscheiden, ob die betreffende
Einrichtung für den Patienten \EE{geeignet} ist und ggfs.
mit dem zuständigen Arzt der Einrichtung Rücksprache zu
halten!\footnote{Eine Bettenzusage erfolgt meist durch die
Einrichtung, ohne dass sie Näheres über das Problem des
Patienten weiß.} Weiters ist es ratsam, vor Transportbeginn
sich die \EE{Bettenzusage telefonisch bestätigen} zu lassen!
}{}{}{}{}


\clearpage
\section{Weiterführendes Patientengespräch und Anamnese: SAMPLER}




\subsection{Strukturierte Anamneseerhebung mittels \enquote*{SAMPLER}-Schema}
\index{SAMPLER}\index{S-KAMEL}

\Layout{27}
{}
{}
{}
{./images/xkcd/desert_island}
{Stille Wasser sind tief \ldots}
{xkcd.com, CC-BY-NC-2.5}

\bigskip

\noindent Als \T{Anamnese} bezeichnet man das Erheben der
Krankengeschichte. Dies schließt sowohl frühere und
chronische Erkrankungen, Operationen und Verletzungen, als
auch auch das Erheben von Details der gegenwärtigen Beschwerden ein. 
Wichtigstes
Instrument zur Erhebung der Anamnese ist das
\E{Patientengespräch}, aber auch eine \I{Fremdanamnese}
(Informationserhebung bei Angehörigen oder anderen
Beteiligten) oder Dokumentationen anderer
Gesundheitseinrichtungen (\E{Patientenbriefe}, \ldots) sind
wichtige Hilfsmittel.




\TextSlide{\TxSAMPLER}{%
In einer Notfallsituation soll das Anamnesegespräch
gleichzeitig alle wichtigen Informationen zu Tage fördern,
andererseits aber effizient und zeitsparend durchgeführt
werden. Um diese Gegensätze miteinander vereinen zu können 
soll es strukturiert durchgeführt werden. Normalerweise
enthält es die folgenden wichtigen Punkte:

\begin{itemize}
    \item \SAMPLER{S} \EE{\color{Base}S}ymptome und Schmerzen
    \item \SAMPLER{A} \EE{\color{Base}A}llergien
    \item \SAMPLER{M} \EE{\color{Base}M}edikamente
    \item \SAMPLER{P} \EE{\color{Base}P}atientengeschichte (Vor- und
    Grunderkrankungen, Operationen, \ldots)
    \item \SAMPLER{L} \EE{\color{Base}L}etzte \ldots
    \item \SAMPLER{E} \EE{\color{Base}E}reignisse vor
    Notfalleintritt
    \item \SAMPLER{R} \EE{\color{Base}R}isikofaktoren
\end{itemize}

In Anlehnung an die Anfangsbuchstaben wird dieses
Abfrageschema \T{SAMPLER}\footnote{Sample \Lang{engl.}:
Probe}\footnote{Manchmal wird dieses Schema 
\I{S-KAMEL} genannt, für Symptome, Krankheiten, Allergien,
Medikamente, Ereignisse und Letztes.} bezeichnet.
Ein derartiges Schema kann im hektischen
Notfall Sicherheit geben. Das Abarbeiten der Fragen
stellt sicher, dass keine wichtigen Informationen
vergessen werden. In der Praxis kann die Reihenfolge der
Fragen je nach Situation abweichen.

}{%
\begin{itemize}
    \item \SAMPLER{S} \EE{\color{Base}S}ymptome und Schmerzen
    \item \SAMPLER{A} \EE{\color{Base}A}llergien
    \item \SAMPLER{M} \EE{\color{Base}M}edikamente
    \item \SAMPLER{P} \EE{\color{Base}P}atientengeschichte 
    \item \SAMPLER{L} \EE{\color{Base}L}etzte \ldots
    \item \SAMPLER{E} \EE{\color{Base}E}reignisse vor
    Notfalleintritt
    \item \SAMPLER{R} \EE{\color{Base}R}isikofaktoren
\end{itemize}
}{}{}{}




\TextSlide{\protect\biolinumGlyph{uni24C8} Symptome und
Schmerzen}{%
\index{Schmerz}Die Einstiegsfrage \Speech{\enquote{Was
kann ich für Sie tun?}} weist oft
bereits auf das Haupt- bzw. Leitsymptom hin. Ist dies
nicht der Fall, so muss als erstes nach dem Hauptsymptom gefragt werden:

\begin{quote}
    \Speech{\enquote{Welche Beschwerden haben Sie?}}
    
    \Speech{\enquote{Was können wir für Sie tun?}}
\end{quote}
Außerdem muss immer auch speziell nach Schmerzen gefragt werden!
\begin{quote}
    \Speech{\enquote{Haben Sie (dabei) Schmerzen?}}
\end{quote}

\subparagraph{Fragen zu Symptomen -- OPQRST}
Gibt der Patient Symptome an, müssen diese weiter abgeklärt
werden, hier am Beispiel Schmerzen:
\index{Schmerz!Arten}

\begin{itemize}
  \item \EE{Lokalisation}: \Speech{\enquote{Wo sind die
  Schmerzen?}}
  \item
  \EE{Qualität}: \Speech{\enquote{Wie ist der Schmerz? Bohrend,
  stechend, brennend, ziehend, krampfartig, auf- und
  abschwellend?}}
  \item \EE{Ausstrahlung}: \Speech{\enquote{Strahlt der Schmerz
  irgendwohin aus?}}
  \item \EE{Zeitdauer / Beginn}: \Speech{\enquote{Seit
  wann haben Sie die Schmerzen? Wie lange haben Sie die Schmerzen schon?
  Haben Sie so einen Schmerz schon einmal gehabt? Haben die
  Schmerzen plötzlich oder langsam eingesetzt?}}
  \item \EE{Geschehnisse zum Zeitpunkt des Auftretens}:
  \Speech{\enquote{Was haben Sie gemacht als es begonnen
  hat?}}
  \item \EE{Provokation / Palliation}:
  \Speech{\enquote{Was macht den Schmerz besser oder schlechter?}}
  \item \EE{Stärke}: \Speech{\enquote{Wie stark ist der
  Schmerz?}}
  Setzen Sie die Aussage des Patienten immer in Relation
  zu seinem Verhalten, seinem Gesichtsausdruck und zu
  seiner Körperhaltung! Einem Patienten mit starken
  Schmerzen sieht man dies oft, aber nicht immer
  an!\ZFootnote{\E{Stärke des Schmerzes}: Einige Autoren
  empfehlen, den Schmerz anhand einer Skala von
  \VonBis{0}{10} beurteilen zu lassen.
  % %   Beurteilung: ex STEUER Review 2011-07-11
  %     Eine solche Beurteilung ist jedoch nur
  %     im Verlauf sinnvoll, da der Patient im akuten
  %     Stadium keinen Anhaltspunkt hat, um die Stärke der
  %     Schmerzen zuverlässig in Relation zu setzen und
  %     einzustufen.
  }
\end{itemize}
Sinngemäß gelten diese Fragen auch für alle anderen
denkbaren Symptome.
\Z{Diese
Fragenabfolge wird in Anlehnung auf die Benennung der
EKG-Zacken als \T{OPQRST} bezeichnet\ZFootnote{\I{OPQRST}: Onset
(Geschehnisse zum Zeitpunkt des Auftretens), Provocation or
Palliation (Was macht es besser bzw. schlechter), Quality (Schmerzqualität,
-art), Region and Radiation (Lokalisation und
Ausstrahlung), Severity (Stärke), Time (zeitlicher
Verlauf)}}.

}{
% \begin{itemize}
%   \item Wo?
%   \item Ausstrahlung?
%   \item Seit wann? Wie lange schon?
%   \item Plötzlicher oder langsamer Beginn?
%   \item Verlauf?
%   \item Wie?
%   \item Wie stark?
%   \item Was provoziert den Schmerz?
% \end{itemize}
% 
\TabHeading{OPQRST}

\begin{itemize}
  \item \EE{Lokalisation}: \Speech{\enquote{Wo sind die
  Schmerzen?}}
  \item
  \EE{Qualität}: \Speech{\enquote{Wie ist der Schmerz? Bohrend,
  stechend, brennend, ziehend, krampfartig, auf- und
  abschwellend?}}
  \item \EE{Ausstrahlung}: \Speech{\enquote{Strahlt der Schmerz
  irgendwohin aus?}}
  \item \EE{Zeitdauer / Beginn}: \Speech{\enquote{Seit
  wann haben Sie die Schmerzen? Wie lange haben Sie die Schmerzen schon?
  Haben Sie so einen Schmerz schon einmal gehabt? Haben die
  Schmerzen plötzlich oder langsam eingesetzt?}}
  \item \EE{Geschehnisse zum Zeitpunkt des Auftretens}:
  \Speech{\enquote{Was haben Sie gemacht als es begonnen
  hat?}}
  \item \EE{Provokation / Palliation}:
  \Speech{\enquote{Was macht den Schmerz besser oder schlechter?}}
  \item \EE{Stärke}: \Speech{\enquote{Wie stark ist der
  Schmerz?}}
\end{itemize}
}{}{}{}





\TextSlide
{Besonderheit: Kolikartiger Schmerz}
{\label{topic:kolik}%
\index{Kolik!Schmerzqualität}%
\index{Schmerz!kolikartiger}%
%
Ein \T{Kolikartiger Schmerz} ist \EE{an-}
und \EE{abschwellend}.
Er wird oft als krampfhaft und sehr stark beschrieben. Er
kommt häufig bei Verstopfung der Gallen- und Harnwege durch
Steine vor (\FullPrettyRef{topic:gallenkolik},
\FullPrettyRef{topic:nierenkolik}).}
{
\begin{compactitem}
  \item Auf- und abschwellend
  \item Krampfhaft
\end{compactitem}
}



\Text{\biolinumGlyph{uni24B6} Allergien}{%
Auf die Frage nach möglichen Allergien wird häufig
vergessen, ein Informationsmangel kann aber fatale
Auswirkungen haben. Es muss daher immer nach Allergien,
in einer zweiten Frage sogar speziell nach
Medikamentenunverträglichkeiten gefragt werden! Bedenke: Ein
Notfallpatient kann bewusstlos werden und vielleicht zu
einem späteren Zeitpunkt nicht mehr befragt werden!
}{}{}{}




\TextSlide{\biolinumGlyph{uni24C2} Medikamente}{%
Die vom Patienten eingenommenen Medikamente lassen
Rückschlüsse auf seine Grunderkrankungen zu und sind im Falle
eines längeren Spitalsaufenthalts eine wertvolle Information.
Auch die Regelmäßigkeit oder Unregelmäßigkeit der Einnahme
kann Anhaltspunkte für das Finden der
Verdachtsdiagnose bzw. der Umstände des Patienten liefern.
Angaben zu Medikamenten sind für die nachbehandelnde
Einrichtung zu dokumentieren. Es muss immer nach
Medikamentenlisten o.\,ä. 
gefragt werden, diese sind für
die weiterbehandelnde Stelle mitzunehmen! Sind keine Listen
vorhanden, so  müssen die entsprechenden
Medikamentenpackungen, Arztbriefe, Rezepte o.\,ä. des
Patienten mitgenommen werden.
Sollte dies vernachlässigt werden, kann es aufgrund der Unkenntnis über die 
Medikation des Patienten zu einer schlechten oder gefährlichen Versorgung 
des Patienten kommen (verspätete Gabe von Medikamenten,
Wechselwirkungen, \ldots).

\Fazit{Immer nach Medikamentenlisten fragen!}

\subparagraph{Fragen zu Medikamenten:}
\begin{quote}
    \emph{\enquote{Welche Medikamente nehmen Sie? Haben
    Sie eine Liste?}}
    
    \emph{\enquote{Wogegen nehmen Sie diese Medikamente?}}
    
    \emph{\enquote{Nehmen Sie Ihre Medikamente regelmäßig?}}
    
    \emph{\enquote{Wann haben Sie Ihre Medikamente zuletzt genommen?}}
    
\end{quote}
}{
\begin{itemize}
  \item Welche?
  \item Wogegen?
  \item Regelmäßige Einnahme?
  \item Wann zuletzt genommen?
  \item \EE{Medikamentenliste}
\end{itemize}
}{}{}{}

\Cave{Das Ignorieren bzw. die Nicht-Mitnahme von 
Medikamentenlisten o.\,ä. stellt einen schweren Fehler dar!}

\Layout{57}{}
{}
{}
{\FontSansNarrowFamily\scriptsize\begin{tabulary}{\linewidth}{LLLLLL}
\toprule\addlinespace
\noindent\TabHeading{(Marken-) Name}
				& \noindent\TabHeading{(Grund-) Erkrankungen}
									& \noindent\TabHeading{Anmerkungen}	
													& \noindent\TabHeading{Andere Namen / Generika}		
																					& \noindent\TabHeading{Wirkstoff} 	
																									& \noindent\TabHeading{Wirkung} 			
\\\addlinespace\midrule\addlinespace
\noindent\Ca{Marcoumar{\texttrademark}}
				& z.\,B. Herzrhytmusstörungen
									& Massiv gesteigerte \EE{Blutungsneigung}
							 						& 								&				& Blutgerinnung ↓ 						
\\\addlinespace
Aspirin{\texttrademark} 		
				& Schmerzen, gripp. Infekte
									& Gesteigerte Blutungsneigung, schädigt Magenschleimhaut
													& \Z{unzählige, z.\,B.: ASS-Genericon{\texttrademark}, ASS-Ratiopharm{\texttrademark}, \ldots }
																					& \Z{Acetyl-Salicylsäure (ASS)} 
																									& Gerinnungshemmend\Z{, schmerzhemmend, entzündungshemmend, fiebersenkend} 
\\\addlinespace
Mexalen{\texttrademark} 		
				& Schmerzen, gripp. Infekte
									& Lebertoxisch, Vergiftungen häufig
													&								& Paracetamol 	& \Z{Fiebersenkend, entzündungshemmend, schmerzhemmend}
																						
\\\addlinespace
Digitalis{\texttrademark} 		
				& Herzinsuffizienz, Herzrhythmusstörungen
									& Über-/Unterdosierungen häufig, kann Herzrhythmusstörungen verursachen
													& 								&				& \Z{Herzkraftsteigernd, Herzrhythmus-stabilisierend}
\\\addlinespace
Diverse; z.\,B.: Voltaren{\texttrademark}, Aspirin{\texttrademark}, \ldots 
				& Schmerzen, gripp. Infekte					
									& Schädigung der Magenschleimhaut
													& Diverse 						& \Z{Nicht-steroidale Antirheumatika (NSAR), z.\,B.: Diclofenac, ASS, \ldots}
																									& \Z{Fiebersenkend, entzündungshemmend, schmerzhemmend} 												
\\\addlinespace
Insulin		& Diabetes mellitus	& Hypoglykämie; Hyperglykämie bei Nichteinnahme				
													& Diverse, verschiedene lang- und kurzwirksame Präparate verfügbar							
																					& Insulin		& Blutzucker ↓												
\\\addlinespace
ThromboASS{\texttrademark}		
			& Z.\,n	Herzinfarkt, Z.\,n. Insult
								&  Anamnese!
													& 							
																					& \Z{Acetyl-Salicylsäure (ASS)}		
																									& \Z{Gerinnungshemmend}											
\\\addlinespace
Plavix{\texttrademark}		
			& Z.\,n	Herzinfarkt, Z.\,n. Insult
								&  Anamnese!
													& 							
																					& \Z{Clopidogrel}		
																									& \Z{Gerinnungshemmend}
\\\addlinespace
Nitroglycerin		
			& Koronare Herzkrankheit
								&  Erhältlich u.\,a. als Spray, Kapseln und Pflaster
													& 							
																					& Nitrolglycerin	
																									& Erweiterung der Herzkrankgefäße\par Achtung: Blutdruck ↓
																					
\\\addlinespace\bottomrule
\end{tabulary}

{\noindent\footnotesize Andere ähnliche Präparate, welche die Gerinnung wie Marcoumar beeinflussen, wären z.\,B. Pradaxa{\texttrademark} oder Sintrom{\texttrademark}}}
{Besonders häufige bzw. wichtige Medikamente}
{}



\TextSlide
{\biolinumGlyph{uni24C5} Patientengeschichte}
{Notfälle ereignen sich manchmal aufgrund der akuten
Verschlechterung einer bestehenden Vorerkrankung. Außerdem
hilft die Beurteilung der Messwerte, wenn man über die
Grunderkrankungen des Patienten Bescheid weiß. Zum Beispiel
kann die Ursache eines hohen Blutdrucks eine bereits
bestehende Erkrankung, aber auch ein Notfallgeschehen sein.
Neben der Befragung können auch Dokumentationen von anderem
Gesundheitsfachpersonal (Spitals- bzw. Arztbriefe, Befunde,
Pflegedokumentation, etc.) Auskunft über bestehende
Vorerkrankungen liefern.

% OP, Dauer-erkr. \ldots

Allgemeine Fragen,
wie z.\,B.
\begin{quote}
    \Speech{\enquote{Leiden Sie unter irgendwelchen
    Krankheiten?}}
    
    \Speech{\enquote{Bestehen bei Ihnen Vorerkrankungen?}}
\end{quote}
können anschließend durch Symptom-spezifische Fragen ergänzt
werden:
\begin{quote}
    \Speech{\enquote{Hatten Sie schon einmal einen
    Herzinfarkt?}}
    
    \Speech{\enquote{Leiden Sie unter hohem Blutdruck?}}
    
    \Speech{\enquote{Sind Sie zuckerkrank?}}
\end{quote}}
{
\begin{itemize}
  \item Grund- und Vorerkrankungen?
  \item Behinderungen?
  \item Operationen?
  \item \ldots
\end{itemize}
}
{}
{}
{}


\TextSlide{\biolinumGlyph{uni24C1} Letzte \ldots}{%
Der letzte Punkt des Abfrageschemas bezieht sich auf letzte
Ereignisse, welche mit dem Notfall in Zusammenhang stehen
könnten. Abhängig von den
Beschwerden, Leitsymptomen bzw. Verdachtsdiagnosen 
können verschiedene Fragen zielführend sein:
\begin{quote}
    \emph{\enquote{Wann haben Sie zuletzt etwas
    gegessen/getrunken?}} (Ist der Patient nüchtern?)
    
    \emph{\enquote{Wann hatten Sie zuletzt Stuhlgang?}}
    (weiterführende Frage: Wie hat der Stuhl ausgesehen?)
    
    \emph{\enquote{Wann war die letzte Regelblutung?}}
    (Besteht die Möglichkeit einer Schwangerschaft?)
    
    \emph{\enquote{Wann waren Sie zuletzt beim Arzt bzw. im
    Spital?}}
\end{quote}
}{
\begin{itemize}
  \item Letzte Mahlzeit?
  \item Letzter Stuhlgang?
  \item Letzte Regelblutung?
  \item Letzter Spitalsaufenthalt?
  \item Letzter Arztbesuch?
  \item Letzte Blutzuckerkontrolle?
  \item \ldots
\end{itemize}
}{}{}{}




\Text{\biolinumGlyph{uni24BA} Ereignisse vor Notfalleintritt}{%
Oft ereignen sich vor einem Notfall Schlüsselereignisse,
welche den Notfall herbeiführen oder diesen auslösen.
Fragen Sie daher immer nach möglichen Besonderheiten oder
Tätigkeiten des Patienten vor dem Notfall:
\begin{quote}
    \Speech{\enquote{Was haben Sie vor Beginn der Symptome
    gemacht?}}
    
    \Speech{\enquote{Haben Sie sich vor Beginn der Symptome
    besonders angestrengt?}}
    
    \ldots{}
\end{quote}
In dieser Kategorie kann es auch hilfreich sein, wenn Sie
Angehörige oder Zeugen befragen:
\begin{itemize}
    \item Hat der Patient vor dem Notfall etwas eingenommen
    (Überdosierung)?
    \item Hat der Patient gekrampft bevor er bewusstlos liegen geblieben ist?
    \item Ist der Patient vor dem Eintreten der Symptome gestürzt?
    \item u.\,v.\,m.
\end{itemize}
}{

}{}{}{}

\TextSlide{\biolinumGlyph{uni24C7} Risikofaktoren}{%
Viele Erkrankungen werden durch das Vorhandensein von
Risikofaktoren begünstigt.
Das Vorliegen von Riskofaktoren bedeutet natürlich nicht,
dass eine bestimme Erkrankung vorliegt, die Kenntnis über
Risikofaktoren kann aber in der Diagnosefindung hilfreich
sein und das Gesamtbild des Patienten vervollständigen.
Beispiele für Riskofaktoren wären:
\begin{itemize}
  \item Alter
  \item Geschlecht
  \item Immobilisierung, Bettlägrigkeit (Thrombosen,
  Lungenembolie, Pneumonie, \ldots)
  \item Rauchen (Lungenkrebs, COPD, Lungenembolie, \ldots)
  \item Diabetes mellitus (Herz-Kreislauf-Erkrankungen:
  Koronare Herzkrankheit, Herzinfarkt, Schlaganfall,
  periphere arterielle Verschlusskrankheit (pAVK),
  Wundheilungsstörungen, \ldots; Niereninsuffizienz, \ldots)
  \item Arterielle Hypertonie (Herz-Kreislauf-Erkrankungen, s.\,o.)
  \item Hormonelle Verhütung (\enquote*{Pille}; Lungenembolie, \ldots)
\end{itemize}
Die Frage nach Risikofaktoren ist abhängig vom Einzelfall,
oft sind wichtige Risikofaktoren auch ohne Nachfragen
erkennbar (Bettlägrigkeit, Alter, Geschlecht \ldots). Es ist
wichtig, sich die gefundenen Riskofaktoren bei der Bewertung
des Patienten zu vergegenwärtigen.
}{
\begin{itemize}
  \item Begünstigen Erkrankungen
  \item Z.\,B.:
  \begin{itemize}
    \item Alter
    \item Geschlecht
    \item Immobilisierung, Bettlägrigkeit
    \item Rauchen
    \item Diabetes mellitus
    \item Arterielle Hypertonie
    \item Hormonelle Verhütung
  \end{itemize}
\end{itemize}
}{}{}{}


\clearpage
\subsection{Dokumentation von anderen Gesundheitseinrichtungen und -berufen}
\label{topic:dokumentation-andere-gesundheitsberufe}

\subsubsection{Arzt- und Entlassungsbriefe}

Nach einer stationären Aufnahme in einer
Gesundheitseinrichtung wird ein Arzt- bzw. Entlassungsbrief
erstellt und dem Patienten übergeben oder zugeschickt. In
diesem Dokument finden sich wichtige Informationen bezüglich
der früheren Krankheitsgeschichte des Patienten, der
Medikation sowie dem Behandlungsverlauf in der jeweiligen
Einrichtung. Eventuell ist auch ein Pflegebericht beigelegt.
Der Brief ist somit eine ideale \E{Ergänzung} des
Anamnesegespräches (aber kann es \E{nicht} ersetzen!).

In jedem Fall ist zu überprüfen, ob es sich tatsächlich um
den Brief des Patienten handelt, und ob die Angaben darin
noch aktuell sind (insbesonders Angaben zu Diagnosen und
Medikation können sehr rasch veralten).

% \bigskip
% \clearpage
\subsubsection{Beispiel: Patient mit schlechtem
Allgemeinzustand}

\hrule

\vspace{1em}

\begin{center}\sffamily
 \textbf{Krankenanstalt der Barmherzigkeit} 
 \\II. Med. Abteilung 
 
Vorstand: Prim. Univ.-Doz. Dr. Vicco von Bülow
 
{\scriptsize Hartlgasse 16a, A-1240 Wien
\\
Tel.: +43\,555\,328\,745\,-\,0 ~~~ Fax: +43\,555\,328\,745\,-\,100}
\end{center}

{%\fontfamily{lmtt}\fontseries{lc}\selectfont                                                                                               
\FontTtNarrowFamily

\vspace{1em}
                                                                                                                            


 Peter Zapfl\\
Diesseitsweg 13a/12/5/34 \\
A-1147 Wien
\vspace{1.5em}

\begin{flushright}
Wien, am 25.07.2009
\end{flushright}

  Wir berichten über den Aufenthalt unseres Patienten Herrn Peter ZAPFL, geb.
  am 13.12.1936, welcher vom 30.03.2009 bis 17.04.2009 an unserer Abteilung
  stationär aufgenommen war.

{\bfseries   Aus der Anamnese:}

\begin{quote}
  Kindheit: \MarginNo{\Z{AE: Appendeketomie, Entfernung des
  Blinddarms; TE: Tonsillektomie, Entfernung der
  Mandeln.}}AE, sowie TE, WK II: längerer Lazarettaufenthalt
  wegen Typhus.
  
  Seit 2003 sind ein bis dato nicht insulinpflichtiger
  Diabetes mellitus, sowie eine art. Hypertonie bekannt.
  
  2005 Aufenthalt an unserer Abteilung wg. eines im
  \MarginNo{\Z{CT:
  Computertomographie.}}CT nachgewiesenen rechtshirnigen
  ischämischen Insultes, anschließend Rehabaufenthalt in Bad
  Schallerbach.
  
  Zuletzt lag der Patient vom 21.11.2008 bis 23.12.2008
  wegen eines nicht insulinpflichtigen, entgleisten Diabetes
  mellitus, bei Z.~n. ischämischen zerebralen Insult 2005
  mit Halbseitenzeichen links und arterieller Hypertonie an
  unserer Abteilung.
  
  Zwischenzeitlich keine Auffälligkeiten.
  
  Alkohol: verneint. Nikotin: verneint. Koffein: fallweise.
  Harn, Stuhl: unauff.
  Schlaf: unauff.
\end{quote}


{\bfseries Bisherige Med.:} 

\begin{quote}
  \begin{tabular}{lll}
Renitec 10 mg &&1-1-1
\\
Diabetex 850 mg &&1-1-1
\\
Diab Diät 12 BE &&2-1-4-2-2-1
\\
Thromboass 100 mg &&1-0-0
\\
\end{tabular}
\end{quote}





 Die nunmehrige Aufnahme erfolgte mit dem ASB, der wegen zunehmender
 Verschlechterung des AZ des Patienten vom Wohnungsnachbarn
 berufen wurde.

{\bfseries Auffälliges im Aufnahmestatus:}

\begin{quote}
 Deutlich herabgesetzter \MarginNo{AZ: Allgemeinzustand\Z{,
 EZ: Ernährungszustand}.}AZ, noch ausreichender EZ, zeitlich
 und örtlich teilorientiert, motorische Unruhe. Aktive
 Beweglichkeit der re. \MarginNo{OE: Obere Extremität, UE:
 Untere Extremität.}OE und UE, beinbetonte spastische
 Halbseitenlähmung links mit gesteigerten Sehnenreflexen
 links bei Z.\,n. re-Hirn-Insult, mäßige
 Artikulationsstörungen.
 
 Haut trocken, \MarginNo{\Z{Turgor: vom Flüssigkeitsgehalt
 abhängiger Spannungszustand der Haut.
 \nocite{Pschy259}}}Turgor herabgesetzt, kein Dekubitus
 nachweisbar.
 \end{quote}

{\bfseries Befunde:}
\begin{quote}
Labor siehe Beilage.

\Margin{\Z{C/P: \enquote{Cor-Pulmo}, Thoraxröntgen.}}C/P.: Keine
wesentliche Änderung gegenüber Vorbefund von 12/2008.

Schädel CT: Keine wesentliche Änderung gegenüber Letztbefund
von 12/2008, Feinbefund siehe siehe Beilage.

Augen: Erstgradige hypertone Veränderungen der
\MarginNo{\Z{Retina: Netzhaut.}}Retinagefässe, kein Hinweis
für \MarginNo{\Z{Diabetische Retinopathie: Schädigung der
Netzhautgefäße im Rahmen des Diabetes mellitus.}}diabetische
Retinopathie.

Neuro: Anamnestisch Z.\,n. Insult 2005 im Strombereich der
A. cerebri ant., mit beinbetonter spastischer Halbseitenlähmung
li und gesteigerten \MarginNo{\Z{BSR, TSR: Sehnenreflexe des
Bizeps und Trizeps.}}BSR und TSR li sowie stark gesteigerten
\MarginNo{\Z{PSR: Patellarsehnenreflex. Babinski-Reflex:
krankhafter Reflex bei Schädigung des Rückenmarks oder
Gehirns.}}PSR li., Babinski li. pos., Trömner u. Knips
beidseits neg., kein Hinweis auf Re-Insult, weiterhin
physi\-ka\-lisch- thera\-peutische und logopädische
Maßnahmen empfohlen. Der durchgeführte Schlucktest ergab
keinen Hinweis auf eine relevante Schluckstörung.

Uro: Prost. rect. dig. unauffällig.

EKG: \MarginNo{\Z{Sinusrhythmus, Herzfrequenz 74,
Linkstyp.}}SR 74, LT, altersentsprechender Kurvenverlauf.

Der Blutdruck war bei mehrfachen Kontrollen gegen Ende des
Aufenthaltes normoton, zuvor im hypotonen Bereich.
\end{quote}




{\bfseries Epikrise:}

\begin{quote}
  Nach Behebung des bei Aufnahme bestehenden
  Flüssigkeitsdefizites mit geeigneter Infusionstherapie,
  besserte sich der AZ des Patienten zusehends. Mit in der
  Folge durchgeführten physi\-ka\-lisch- therap. Maßnahmen
  konnte der Mobilisisationsgrad des Pat. günstig
  beeinflusst werden.
  
  Der Pat. wird bis zum selbständigen Gehen mit Rollator
  mobilisiert entlassen. Die zu Aufenthaltsbeginn durchwegs
  im hypotonen Bereich gemessenen RR-Werte haben sich in
  mehrfachen Kontrollen unter dem angegeben antihypertonen
  Regime normalisiert. Ebenso hat sich die ursprünglich
  unbefriedigend vorgefundene diabetische Stoffwechsellage,
  bei konsequenter Einhaltung des vorgeschlagenen Regimes im
  therapeutischen Bereich stabilisiert.
  
  Im Einverständnis mit dem Pat. und dessen Tochter wurde
  ein Platz im Bernhard-Prischl-Heim zur Gewährleistung
  einer regelmäßigen und auch ausreichenden
  Flüssigkeitszufuhr, sowie einer gesicherten Diabetes-
  Diät, organisiert. Eine Kontrolle der vor etwa einem Jahr
  durchgeführten \MarginNo{\Z{Carotissonographie:
  Ultraschall der Halsschlagader.}}Carotissonographie ist
  angezeigt. Ein entspr. Termin wurde am 27.12.2009  12,30
  Uhr an der Angiolog Ambulanz d. KA Mitschka-Stiftung
  fixiert. Eine Befund\-be\-sprech\-ung ist nach
  Terminvereinbarung an h.o. Int. Ambulanz möglich.
\end{quote}

\MarginNo{\Z{}}
{\bfseries Hauptdiagnose:}

\begin{quote}
\begin{tabular}{ll}
Hypohydratation & E86.0
\\
\end{tabular}
\end{quote}

{\bfseries Nebendiagnosen:}

\begin{quote}
\begin{tabular}{ll}
entgleister. Diabetes mellitus. (NIDDM) & E12.31, E12.40,
E12.72
\\
Art. Hypertonie & I10.1
\\
Z.n. Zerebralem Insult mit HSS li 2005& I63.3
\\
Spastische Hemiparese li, beinbetont&I69.3, G81.1
\\
\end{tabular}
\end{quote}



{\bfseries Therapie:}

\begin{quote}
\begin{tabular}{lll}
Renistad Tab. 5,0 mg && 1-1-1 
\\
Glucophage 850 mg Tab && 1-1-1 
\\
Diab. Diät 12 BE&&(2-1-4-2-2-1) 
\\
Thromboass Tab.&&1-0-0 
\\
tgl. Flktszufuhr nicht unter 2 l&&
\end{tabular}
\end{quote}




{\bfseries Kontrollen:}

\begin{quote}
  NFP, Harnbefund, Elektrolyte, RR, BZ, (Patient. führt seit
  Jahren Selbstkontrollen durch und führt darüber
  Aufzeichnungen)) weiterhin neurolog. und ophthalmolog.
  Kontrollen.
\end{quote}









Mit kollegialen Grüßen\vspace{2em}\\ Dr. Adolf Schrammel, Stationsarzt

\Hand{Schrammel}

vidit OA Univ.Doz. Dr. Alois Schremser
}

\vspace{2em}

\hrule




% \clearpage
% \subsubsection{Beispiel: Patient nach Operation}
% 
% \hrule
% 
% \vspace{1em}
% 
% \begin{center}\sffamily
%  \textbf{Landesklinikum Hinter a. d. Tupfing}
%  \\Abteilung für Chirurgie
% 
% Vorstand: Prim. Dr. Viktor Hauer
% 
% {\scriptsize Spitalgasse 2, A-3212 Hinter/Tupfing
% \\
% Tel.: +43\,555\,328\,745\,-\,0 ~~~ Fax: +43\,555\,328\,745\,-\,100}
% \end{center}
% 
% {%\fontfamily{lmtt}\fontseries{lc}\selectfont
% \sffamily
% 
% \vspace{1em}
% 
% 
% 
%  Walter Mayer\\
% Tako-Tsubo-Weg 14a/11/25 \\
% A-1147 Wien
% \vspace{1.5em}
% 
% \begin{flushright}
% Hinter/Tupfing, am 25.08.2012
% \end{flushright}
% 
%   Wir berichten über den Aufenthalt unseres Patienten Herrn Walter M\,A\,Y\,E\,R, geb.
%   am 01.01.1933, welcher vom 23.08.2012 bis 25.08.2012 an
%   unserer Abteilung stationär aufgenommen war.
% 
% 
% \begin{description}
%   \item[Grund der Aufnahme:] Obstipation, V. a. Ileus
%   \item[Wichtiges zur Aufnahme und Anamnese:] Herr Mayer war
%   bereits wiederholt h.\,o. wegen Obstipation bei bekanntem
%   N. recti in Behandlung. Eine Rektumamputation mit Anus
%   praeter wurde vom Patienten bisher abgelehnt. Zuletzt erfolgte eine palliative
%   Tumorabtragung. Numehr
%   erfolgt die stationäre Aufnahme über die zentrale
%   Notfallambulanz wegen Stuhlverhalt, angeblich seit 2 Wochen.
% \end{description}
% 
% 
% \begin{description}
%   \item[Status präsens:] Reduzierter AZ, reduzierter EZ, \\
%   Cor: rhythmisch, tachykard, rein; Pulmo: vesikuläres
%   Atemgeräusch beidseits, sonorer Klopfschall, Sprechdyspnoe
%   Abdomen: weich, spärliche Darmgeräusche, kein
%   Druckschmerz, Stuhlwalzen\\
%   Neurologisch; grob neurologisch unauffällig
%   \item[Auffälliges im Aufnahmelabor:] siehe Beilage
%   \item[EKG:] VH-Flimmern, 80/min, LAHB
%   \item[Cor-Pulmo-Röntgen:] Keine Infiltrate, keine
%   Stauungszeichen
%   \item[Digital-rektale Untersuchung:] Resistenz ca. 4 cm ab
%   ano, Ampulle stuhlfrei
% \end{description}
% 
% \paragraph{Verlauf:}
% Der Patient wird umgehend auf der Herzüberwachungsstation
% aufgenommen und überwacht.
% 
% \paragraph{Diagnosen}~
% \par
% 
% \begin{tabulary}{\linewidth}{rL}
% --.-
% &
% \textbf{Tachykardes Vorhofflattern}
% \\\addlinespace
% & Z.n. Vorhofflimmern und medikamentöser Cardioversion
% 12/2011
% \\
% &Z.n. Pulmonalvenenisolation 3/2012
% \\
% &Koronarsklerose
% \\
% & Ausschluss sign. Stenose 3/2012
% \\
% I10.1
% &
% Art. Hypertonie
% \\
% &Z.n. AE, TE
% \end{tabulary}
% 
% \paragraph{Therapie bei Entlassung}~
% 
% \begin{tabulary}{\linewidth}{Ll}
% Sedacoron 200mg & 1-0-0
% \\
% Concor 5mg & 1-0-0
% \\
% ThromboASS 100mg & 1-0-0
% \\
% Simvastatin 20mg & 0-0-1
% \\
% Marcoumar & lt. Pass
% \end{tabulary}
% 
% 
% \paragraph{Weiteres Procedere:} ~
% 
% 
% 
% 
% \bigbreak
% 
% 
% 
% \noindent\begin{tabularx}{\linewidth}{XXX}
% Oberarzt & & Sekundararzt
% \\
% \textbf{OA Dr. R. Glocke}
% &
% &
% Sek.-Dr. H. Harrer
% \\\addlinespace
% \\
% \Hand{Glocke}
% 
% \end{tabularx}
% \vspace{2em}
% 
% \hrule
% }





% 
% 
% 
% 
% 
% 
% 
% 
% 
% 


% \clearpage
% \subsubsection{Beispiel: Patient mit tachykarder
% Rhythmusstörung}
% 
% \hrule
% 
% \vspace{1em}
% 
% \begin{center}\sffamily
%  \textbf{Landesklinikum Hinter a. d. Tupfing}
%  \\Abteilung für interne Medizin
% 
% Vorstand: Prim. Dr. Viktor Hauer
% 
% {\scriptsize Spitalgasse 2, A-3212 Hinter/Tupfing
% \\
% Tel.: +43\,555\,328\,745\,-\,0 ~~~ Fax: +43\,555\,328\,745\,-\,100}
% \end{center}
% 
% {%\fontfamily{lmtt}\fontseries{lc}\selectfont
% \sffamily
% 
% \vspace{1em}
% 
% 
% 
%  Walter Mayer\\
% Tako-Tsubo-Weg 14a/11/25 \\
% A-1147 Wien
% \vspace{1.5em}
% 
% \begin{flushright}
% Hinter/Tupfing, am 25.08.2012
% \end{flushright}
% 
%   Wir berichten über den Aufenthalt unseres Patienten Herrn Walter M\,A\,Y\,E\,R, geb.
%   am 01.01.1970, welcher vom 23.08.2012 bis 25.08.2012 an
%   unserer Abteilung stationär aufgenommen war.
% 
% 
% \begin{description}
%   \item[Grund der Aufnahme:] Tachykarde Rhythmusstörung,
%   Schmalkomplextachykardie
%   \item[Wichtiges zur Aufnahme und Anamnese:] Herr Mayer war
%   wiederholt h.\,o. nach Episoden von Herzklopfen und Z.\,n. Kollaps vorstellig. Im Dezember
%   2011 konnte erstmalig ein Vorhofflimmern
%   im EKG dargestellt werden. Seitdem besteht eine orale
%   Antokoagulation mit Marcoumar und eine Rezidivprophylaxe
%   mit Sedacoron und Concor. Eine im LK Krems durchgeführt
%   Koronarangiographie ergab keine mäiggradige
%   Koronarsklerose ohne Hinweis auf eine signifikante
%   Stenosierung. Es wurde ferner im März 2012 eine
%   Pulmonalvenenisolation durchgeführt.
% \end{description}
% 
%   \paragraph{Grund der Aufnahme:} Tachykarde
%   Rhythmusstörung, Schmalkomplextachykardie
%   \paragraph{Wichtiges zur Aufnahme und Anamnese:} Herr
%   Mayer war wiederholt h.\,o. nach Episoden von Herzklopfen und Z.\,n. Kollaps vorstellig. Im Dezember
%   2011 konnte erstmalig ein Vorhofflimmern
%   im EKG dargestellt werden. Seitdem besteht eine orale
%   Antokoagulation mit Marcoumar und eine Rezidivprophylaxe
%   mit Sedacoron und Concor. Eine im LK Krems durchgeführt
%   Koronarangiographie ergab keine mäiggradige
%   Koronarsklerose ohne Hinweis auf eine signifikante
%   Stenosierung. Es wurde ferner im März 2012 eine
%   Pulmonalvenenisolation durchgeführt.
% 
%   Der Patient kommt nun bei Z.\,n. Kollaps und
%   Palpitationen. Im Aufnahme-EKG zeigt sich ein tachykardes
%   Vorhofflattern mit 2:1-Überleitung.
% 
% \begin{description}
%   \item[Status präsens:] Reduzierter AZ, adipöser EZ, \\
%   Cor: rhythmisch, tachykard, rein; Pulmo: vesikuläres
%   Atemgeräusch beidseits, sonorer Klopfschall, Sprechdyspnoe
%   Abdomen: weich, unauff. Darmgeräusche, kein Druckschmerz\\
%   Neurologisch; grob neurologisch unauffällig
%   \item[Auffälliges im Aufnahmelabor:] siehe Beilage
%   \item[EKG:] VH-Flattern, 140/min, LAHB
%   \item[Cor-Pulmo-Röntgen:] Keine Infiltrate, keine
%   Stauungszeichen
%   \item[]
% \end{description}
% 
% \paragraph{Verlauf:}
% Der Patient wird umgehend auf der Herzüberwachungsstation
% aufgenommen und überwacht.
% 
% \paragraph{Diagnosen}~
% \par
% 
% \begin{tabulary}{\linewidth}{rL}
% --.-
% &
% \textbf{Tachykardes Vorhofflattern}
% \\\addlinespace
% & Z.n. Vorhofflimmern und medikamentöser Cardioversion
% 12/2011
% \\
% &Z.n. Pulmonalvenenisolation 3/2012
% \\
% &Koronarsklerose
% \\
% & Ausschluss sign. Stenose 3/2012
% \\
% I10.1
% &
% Art. Hypertonie
% \\
% &Z.n. AE, TE
% \end{tabulary}
% 
% \paragraph{Therapie bei Entlassung}~
% 
% \begin{tabulary}{\linewidth}{Ll}
% Sedacoron 200mg & 1-0-0
% \\
% Concor 5mg & 1-0-0
% \\
% ThromboASS 100mg & 1-0-0
% \\
% Simvastatin 20mg & 0-0-1
% \\
% Marcoumar & lt. Pass
% \end{tabulary}
% 
% 
% \paragraph{Weiteres Procedere:} ~
% 
% 
% 
% 
% \bigbreak
% 
% 
% 
% \noindent\begin{tabularx}{\linewidth}{XXX}
% Oberarzt & & Sekundararzt
% \\
% \textbf{OA Dr. R. Glocke}
% &
% &
% Sek.-Dr. H. Harrer
% \\\addlinespace
% \\
% \Hand{Glocke}
% 
% \end{tabularx}
% \vspace{2em}
% 
% \hrule
% }

\begin{Literary}
\fullcite{DahmerGespraechsfuehrung5}

\fullcite{FamPropAeGeFue2010}

\fullcite{ThunRedenI46}

\fullcite{ArgelanderErstinterview8}
\end{Literary}



\section{Weiterführende Diagnostik und
Maßnahmen}


Die Kurzübersicht in 
\FullPrettyRef{tab:checkliste-patientenmanagement}, gibt einen
Überblick über die mögliche weiterführende Diagnostik, die
Anamnese und Maßnahmen, wobei die
Reihenfolge bei jedem Patienten anders ist!
Unter \FullPrettyRef{m:standardmassnahmen-immer}, 
 sind \EE{immer
durchzuführende Standardmaßnahmen} definiert, diese 
müssen bei jeder Patientenbetreuung -- sofern angemessen --
durchgeführt werden.



% M %%%%%%%%%%%%%%%%%%%%%%%%%%%%%%%%%%%%%%%%%%%%%%%%%%% M %
% Immer durchzuführende Standardmaßnahmen
\ProcedurePrintDefault{YY11000C}

% \subsection{Der Patient endet nicht mit der Übergabe}
% 
% \begin{quote}
% Aortenaneurysma
% \end{quote}
% 
% \begin{quote}
% Aortenaneurysma 2
% \end{quote}
% 
% \subsection{Fehler und Fallen}
% vdfvrwbwrtbwrtbrwtb
% 
% \begin{quote}
% Ketanest --- {\itshape Eines Tages, es war Winter, die Straßen verschneit und
% glatt, wurden wir zu einer Patientin mit einer offenen Unterschenkelfraktur
% berufen. Die Dame war auf einer nassen Stufe im Hauseingang augerutscht. Sie
% wirkte soweit gefasst, aufgrund des Schocks\footnote{Schock: hier
% psychologischer Schock.} schien sie noch keine großen Schmerzen wahrzunehemn.
% Trotzdem wurde zur prophylaktischen Schmerztherapie ein Notarzt nachbestellt.
% Dieser entschloß sich zu einer Sedierung, u.\,a. mit
% Dormicum{\texttrademark}\footnote{Dormicum{\texttrademark}, Wikstoff Midazolam,
% ist ein Schlafmittel, welches zur Sedierung und Narkose verwendet wird.}. (An
% die anderen Medikemnte kann ich mich nicht mehr genau erinnern, ich denke
% Ketanest{\texttrademark}\footnote{Ketanest{\texttrademark}, Wikstoff S-Ketamin,
% ist ein Narkosemittel.} war auch mit ihm Spiel.) 
% 
% Ich erninnere mich noch dass der Notarzt etwas vom Dormicum{\texttrademark}
% spritzte und sich nach einiger Zeit wunderte weildie erwünschte Wirkung
% ausblieb. Er spritzte noch etwas nach und langsam begann es zu wirken \ldots 
% \emph{\enquote{Na bumm, die hat aber viel gebraucht}}, merkte er nach der Übergabe
% dann  noch an. Des Rätsels Lösung war aber nicht bei der Patientin zu suchen,
% sondern bei den verwendeten Ampullen: 
% 
% Wir vewendeten Ampullen mit 5\,mg Wirkstoff, die Organisation, der der Notarzt
% angehörte, Ampullen mit 15\,mg (!) Wirkstoff. Daher war es wenig verwunderlich,
% dass \enquote{unser} Dormicum{\texttrademark} viel weniger Wirkung hatte. Nur gut
% dass der Notarzt \enquote{sein} Ketanest{\texttrademark} aus der Notarzt-Tasche
% verwendet hatte: Unsere Ketanest{\texttrademark}-Ampullen hatten einen Inhalt von 50\,mg, die des
% Notarztes 25\,mg {\ldots} }
% \end{quote}







% \clearpage % TODO Temp clearpage


% \clearpage % TODO temporary clearpage
\section{Übergabe an weiterbehandelndes Fachpersonal}  

\Layout{61}{}{%
Die Übergabe des Patienten in den Verantwortungsbereich
eines anderen Sanitäters, eines Notarztes oder des
Spitalspersonals hat geordnet und gewissenhaft zu erfolgen,
damit keine wichtigen Informationen verloren gehen. Diese
\E{Schnittstelle} zwischen Sanitäter und Arzt bzw.
Pflegepersonal kann auf den weiteren Behandlungsverlauf
einen großen Einfluss nehmen und darf daher in seiner
Bedeutung nicht unterschätzt werden! Es empfiehlt sich --
soweit Zeit dafür ist -- sich bereits vor Eintreffen des
Notarztes oder vor Ankunft im Krankenhaus zu überlegen, welche
Informationen wesentlich sind und daher auf keinen Fall
vergessen werden dürfen.

\bigskip\noindent Das Übergabegespräch muss mindestens
\E{folgende Informationen} beinhalten:
\begin{enumerate}
    \item Name, Alter und Geschlecht des
    Patienten
    \item Grund der Berufung, Hauptbeschwerde
    \item Beschreibung der Symptomatik und deren Verlauf
    \item Wesentliche Befunde und Teile der Anamnese, welche
    \E{direkten Bezug zur Beschwerde} haben
    \item Durchgeführte Maßnahmen
    \item Weitere Befunde 
    \item Weitere Anamnese
    \item Sonstige, wichtige Hinweise (z.B. Infektionsgefahren,
    Bluterkrankheit, Allergien etc.)
\end{enumerate}
}
{}
{./images/shared/aufnahmearzt-ihres-vertrauens-edited2-0500.jpg}
{\enquote{Sie wünschen?}}
{Gabriel}

\subsection{Anwendung in der Praxis}

\Text{Primäreinsatz}{%
Das hier vorgestellte Schema wurde primär in Hinblick auf
Primäreinsätze entwickelt und ist dementsprechend anwendbar.
}{}{}{}{}

\Text{Sekundäreinsatz}{%
Für den Sekundäreinsatz (insbesonders Transferierungen von
behandlungspflichtigen Patienten) gilt grundsätzlich das
gleiche Vorgehen wie beim Primäreinsatz. Die Informationen
der abgebenden Stelle bezüglich des Patienten und des
Transports müssen besonders berücksichtigt werden. Die
einzelnen Handlungen müssen sinngemäß an die Situation
angepasst werden.

Auch im Krankentransportdienst (inkl. Heimtransporte) soll
der Einschätzungsblock bis zu Punkt \EE{C} in
angemessener Weise durchgeführt werden. In der Praxis ist
das binnen weniger Sekunden möglich. 
An dieser Stelle sei nochmals auf die \EE{besondere
Gefährdung von \E{Dialysepatienten}} durch
Elektrolytstörungen hingewiesen!
}{}{}{}{}


\Text{Ambulanzdienst}{%
Für Ambulanzdienste gilt grundsätzlich das gleiche Vorgehen
wie beim Primäreinsatz, die einzelnen Handlungen müssen
sinngemäß an die Situation angepasst werden. Auch hier soll
der Einschätzungsblock bis zu Punkt \EE{C} in
angemessener Weise durchgeführt werden.
}{}{}{}{}

%=========Ende Abschnitt Emhofer==========================================================



\begin{Literary}
   \fullcite{AMLS3}

   \fullcite{emergency9}

   \fullcite{ITLSde6}

   \fullcite{PHTLS6}

   \fullcite{SemmelABC1}

   \fullcite{SudoweProfHandelnRD1}

   \fullcite{KirschnickPflege2}

   \fullcite{RueckerBildatlas1}

   \fullcite{Ziesenitz2009}
\end{Literary}

% \clearpage
% \section{Fallbeispiele}
% \subsection{Primäreinsatz}
% \subsubsection{Ein Beispiel: Patient mit Atemnot}
% \Margin{\includegraphics[width=\linewidth]{./images/teaser/rtw-asb-921-alt.jpg}}
% \begin{verbatim}
% Zeit: dd.04.2003 14:15
% Code: 26-B-1 
%    Kranke Person (Differentialdiagnose): 
%    Unbekannter Zustand (Anrufer 3. Hand) 
% BO  : 1210, xxxstraße 13/8/8/45, WHG
% 
% \end{verbatim}
% 
% \begin{quote}
%   \marginlabel{Szeneüberblick}Die Anfahrt gestaltet sich problemlos. Der
%   Berufungsort befindet sich in einem Wohnblock im achten Stockwerk. Das
%   Eingangstor wird fixiert um evtl. nachzufordernden Einsatzkräften (Notarzt,
%   Polizei, Feuerwehr, \ldots ) den Zugang zu erleichtern. Der Aufzug
%   funktioniert und scheint groß genug für die Trage zu sein. 
%   
%   Sie kommen in eine Wohnung. Die Umgebung erscheint sicher, Sie erkennen auch
%   keine Hinweiszeichen auf ein Haustier. Die Wohnung sieht sauber und ordentlich
%   aus. Auf der Couch liegt flach ein \marginlabel{Ersteindruck}50-jähriger Mann mit weit aufgerissenen Augen und
%   sichtlicher Atemnot. Er wirkt blass und schweißig und fasst sich mit der
%   rechten Hand an den Brustbereich.
%   
%   Sie begrüßen den Patienten, stellen sich vor und
%   \marginlabel{Hauptbeschwerde}fragen welche Beschwerden er hat. Der Patient
%   \marginlabel{Atemweg, Atmung}atmet schnell und ringt nach Luft, besondere
%   Atemgeräusche sind nicht hörbar. Er klagt über Schmerzen in der Brust. Diese
%   haben vor einer Viertelstunde angefangen. \marginlabel{Kreislauf}Währenddessen
%   haben Sie an der Radialarterie des Patienten den Puls getastet und eine
%   Rekap-Probe gemacht. Der Puls ist sehr schwach, schnell und unregelmäßig. Die
%   Rekap-Zeit war etwas verlangsamt. Die Haut ist kühl und
%   schweißig.
% \end{quote}
% 
% \noindent Welche Informationen haben Sie bis jetzt gewonnen?
% \begin{quote}
%   Der Patient ist bei Bewusstsein. Die Atemwege scheinen frei, aber er hat
%   Atemnot und einen plötzlich aufgetretenen Thoraxschmerz. Er ist blass und
%   schweißig, vielleicht steckt ein Schock dahinter. Der schwach fühlbare
%   Puls und die verlängerte Rekap-Zeit unterstützen diese Vermutung.
% \end{quote}
% 
% \noindent Ist der Patient vital bedroht?
% \begin{quote}
%   \marginlabel{Vitale Bedrohung: frühzeitig erkannt (noch
%   vor Abschluss des Einschätzungsblockes)}Ja. Zwar wurden
%   noch keine Vitalwerte erhoben, aber bereits jetzt gilt der
%   Patient als vital bedroht aufgrund der plötzlich
%   einsetzenden Thoraxschmerzen und der Atemnot. Außerdem
%   lässt die Kaltschweißigkeit, der schwache Puls und die
%   verlängerte Rekap-Zeit auch nichts gutes vermuten.
% \end{quote}
% 
% \noindent Können Sie schon etwas tun?
% 
% \begin{quote}
%   \marginlabel{\TxMassVit}Ja.  Bei erhaltenem Bewusstsein, Atemnot und
%   Thoraxschmerz wird der Patient mit erhöhtem Oberkörper gelagert. Weiters sprechen
%   die Symptome für den Einsatz von Sauerstoff. Es muss noch schnell geklärt werden
%   ob eine Hyperventilationstetanie (nicht wahrscheinlich) oder COPD-Erkrankung
%   (erfragen) vorliegt.
%   
%   \TxMassVit:
%   
%   \begin{itemize}
%     \item Der Oberkörper des Patienten wird hochgelagert (Atemnot,
%     Thoraxschmerz)
%     \item Der Notarzt wird nachgefordert (Berufungsgrund: akuter
%     Thoraxschmerz mit Atemnot und schlechtem AZ).
%     \item Der Patient wird gefragt ob er an COPD leidet, dies wird verneint.
%     Eine psychisch bedingte Hyperventilation erscheint unwahrscheinlich.
%     
%     Daher wird eine Sauerstoffmaske mit einem Flow von
%     10{\LPerMin*} angelegt.
%     \item Reanimationsbereitschaft wird hergestellt
%     \item Monitoring: Der Patient wird beobachtet, das
%     Fachpersonal nimmt sich vor, sobald Zeit ist, ein
%     Pulsoxymeter und später ein EKG-Gerät anzuhängen.
%   \end{itemize}
%   
% \end{quote}
% 
% 
% \noindent Was sind die nächsten Schritte?
% 
% \begin{quote}
%   \marginlabel{D: Schnelle Untersuchung}Wichtig sind
%   nun die Vitalwerte und das frühe Beginnen des Monitorings.
%   Dann muss man genaueres über das Beschwerdebild in Erfahrung bringen. Der weitere Plan hängt davon ab.
%   
%   Die Vitalwerte werden gemessen und ein Pulsoxy angehängt (RR
%   \RRmmHg{100}{60}, HF 110\PerMin*, Sp\ce{O2}
%   94\PC*). Der Patient gibt an, dass er stärkste
%   stechende/brennende Schmerzen hinter dem Brustbein hat die ins Kiefer ausstrahlen. Außerdem fühle es sich so an als würde jemand auf dem
%   Patienten drauf stehen. Die Schmerzen sind nicht atem- oder
%   bewegungsabhängig  und haben vor einer Viertelstunde
%   {\quotedblbase}wie aus dem Nichts heraus{\textquotedblleft} begonnen.
%   Außerdem bekomme er schlecht Luft, sonst hätte er keine Beschwerden.
%   
%   Der Patient gibt entsprechende Antworten auf Fragen zu
%   seiner Person, der Situation, Zeit und Ort, ist also
%   orientiert. Die Pupillen sind mittelweit und reagieren
%   prompt und seitengleich auf Licht. Die Kraft der
%   oberen Extremität ist schwach, vermutlich aufgrund der
%   belastenden Situation, aber seitengleich, es gibt keinen
%   Hinweis auf neurologische Ausfälle.
% \end{quote}
% 
% \marginlabel{Problem des Patienten}Als Leitsymptome erkennen Sie den
% Thoraxschmerz und die Atemnot.
% 
% \noindent An welche Differentialdiagnosen müssen Sie nun
% denken?
% 
% % \TableWide
% % 	{Titel}%1
% % 	{Beschreibung}%2
% % 	{Label}%3
% % 	{Quelle}%4
% % 	{Lizenz}%5
% % 	{Tabelle}%6
% % 	{Float}%8
% \TableWide
% {Thoraxschmerz: Differentialdiagnosen des
% Patienten}%1
% {Eine Übersicht zum Fallbeispiel}%2
% {}%3
% {}%4
% {}%5
% {%
% \begin{tabularx}{\linewidth}[H]{XXXc}
% \addlinespace\toprule\addlinespace
% \multicolumn{4}{c}{\centering\TabHeading{} }
% \\\addlinespace\midrule\addlinespace
% \centering \TabSub{Verdachtsdiagnose}
%  &
% \centering \TabSub{dafür spricht}
%  &
% \centering \TabSub{dagegen spricht}
%  &
% \TabSub{Bewertung}
% \\\addlinespace\cmidrule(lr){1-1}\cmidrule(lr){2-2}\cmidrule(lr){3-3}\cmidrule(lr){4-4}\addlinespace
% \EE{Herzinfarkt, Angina pectoris}
%  &
% Schmerzen hinter dem Brustbein, Ausstrahlung, nicht atemabhängig,
% Druck-/Engegefühl, Atemnot
%  &
%  &
% sehr wahrscheinlich
% \\\addlinespace
% \EE{Lungenembolie}
%  &
% Schmerzen hinter dem Brustbein, Atemnot
%  &
% untypische Ausstrahlung, 
%  &
% gut möglich
% \\\addlinespace
% \EE{Pneumothorax}
%  &
% Schmerzen im Brustbereich, Atemnot
%  &
% plötzliche stärkste Schmerzen
%  &
% eher unwahrscheinlich
% \\\addlinespace
% \EE{Trauma}
%  &
% Schmerzen
%  &
% kein Hinweis
%  &
% unwahrscheinlich
% \\\addlinespace
% \EE{Erkrankung des Stütz- und Bewegungsapparats}
%  &
% Schmerzen 
%  &
% untypische Ausstrahlung, Atemnot, nicht bewegungsabhängig
%  &
% sehr unwahrscheinlich
% \\\addlinespace\bottomrule
% \end{tabularx}
% }%6
% {H}%8
% 
% 
% Die wahrscheinlichsten Diagnosen sind demnach ein
% Herzinfarkt oder eine Lungenembolie.
% 
% \noindent Fallen Ihnen bei diesen Erkrankungen wichtige
% weitere Untersuchungen ein?
% 
% \begin{quote}
%   \marginlabel{Durchatmen, nachdenken, planen}Bei beiden Erkrankungen ist es
%   wichtig beide Beine zu begutachten. Im Falle einer Lungenembolie sucht man nach
%   einer Beinvenenthrombose (einseitig überwärmtes, bläulich-rötliches geschwollenes
%   Bein). Bei einer Herzinsuffizienz bzw. Herzinfarkt beurteilt man den Rückstau des
%   Blutes, dh. ob die Beine geschwollen sind.
% \end{quote}
% 
% 
% \noindent Außerdem, hat der Patient so etwas schon einmal
% gehabt? Oder sind einschlägige Erkrankungen bekannt?  
% 
% \begin{quote}
%   Sie begutachten beide Beine des Patienten und stellen mittels
%   Daumendruckprobe fest, dass beide Beine leicht geschwollen sind. Auf
%   Nachfrage gibt der Patient an, dass er sonst keine geschwollenen Beine
%   hat.
% \end{quote}
% 
% \noindent Was sind die nächsten Schritte?
% 
% \begin{quote}
%   Die wichtigsten Untersuchungen wurden gemacht, nun muss man sich um die
%   Anamnese kümmern.
%   
%   Symptome \& Schmerzen: siehe oben.
%   
%   Allergien: keine bekannt.
%   
%   Medikamente: Glucophage 850\Mg* \enquote{für den Zucker},
%   Tritace \enquote{für den Blutdruck}. Medikamente wurden heute wie verschrieben eingenommen.
%   
%   Krankheiten: Der Patient leidet an nicht-Insulinpflichtigem Diabetes
%   mellitus und Bluthochdruck. Sonst sind keine Krankheiten bekannt.
%   Spitalsbriefe gibt es nicht.
%   
%   Sie merken sich, dass Sie nachher den Blutzucker messen wollen!
%   
%   Letzte(r):
%   
%   Spitalsaufenthalt: Vor 10 Jahren in der Rudolfstiftung wegen Rotlauf.
%   
%   Arztbesuch: Vor einer Woche wegen Neuausstellung eines Rezeptes.
%   
%   Mahlzeit: Mittagessen, Gulasch
%   
%   Ereignisse vor Notfalleintritt: Schmerzen begannen plötzlich als der
%   Patient Wäsche waschen wollte. In den letzten Wochen immer wieder
%   Stiche in der Herzgegend, diese hat er aber nicht ernst genommen.
%     
% \end{quote}
% % pro:
% % 
% % Lungenembolie 
% 
% \noindent Was glauben Sie wie es weiter geht?
% 
% % \TakeNotesLarge
% \bigskip

% ========Subsection eingefügt vom Emhofer 1.5.2009===============================





\addtocategory{PAM}{AMLS3,emergency9,ITLSde6,PHTLS6,SemmelABC1,SudoweProfHandelnRD1,
KirschnickPflege2,RueckerBildatlas1,Ziesenitz2009}
\addtocategory{BestOf}{ITLSde6,PHTLS6,RueckerBildatlas1}
\ChapterEnd
